\documentclass[11pt,a5paper,twoside]{book}

\usepackage{inputenc}
\usepackage{fontenc}
\usepackage{graphicx}
\usepackage{amsmath}
\usepackage{xcolor}
\usepackage{setspace}
\usepackage{titlesec}
\usepackage{tcolorbox}
\usepackage{colortbl}
\usepackage{helvet}
\usepackage[top=1.5cm,bottom=1.5cm,left=1.5cm,right=1.5cm]{geometry}
\usepackage[spanish]{babel}

\setcounter{secnumdepth}{1}
\setstretch{1.25}

\definecolor{chaptercolor}{HTML}{24527A}
\definecolor{primary}{HTML}{24527A}
\definecolor{secondary}{HTML}{4F81A1}
\definecolor{codebg}{HTML}{F4F6F8}

\titleformat{\chapter}
  {\Huge\bfseries\color{chaptercolor}}
  {\thechapter}
  {1em}
  {}

\begin{document}


\setcounter{tocdepth}{1}
\tableofcontents
\chapter{Fundamentos de Python}

\section{Introducción a Python}
\subsection{Introducción}
Python es un \textbf{lenguaje de programación de propósito general}, lo que significa que puede utilizarse para desarrollar aplicaciones muy diferentes: desde pequeños scripts hasta aplicaciones web, programas de escritorio, herramientas de automatización, inteligencia artificial o análisis de datos.
Una de las principales características de Python es su \textbf{simplicidad}. Su sintaxis está diseñada para ser clara y fácil de leer, lo que permite centrarse en resolver problemas en lugar de preocuparse por los detalles del lenguaje.
Gracias a ello, Python se ha convertido en uno de los lenguajes de programación más populares del mundo y es utilizado tanto por personas que están aprendiendo a programar como por desarrolladores profesionales y empresas de todos los tamaños.
En sus comienzos, Python se utilizaba principalmente en tareas de automatización, prototipado y desarrollo de herramientas, ya que el hardware de la época hacía que otros lenguajes fueran preferibles cuando el rendimiento era crítico. Sin embargo, el aumento de la capacidad de procesamiento de los equipos y la enorme evolución de su ecosistema de bibliotecas han convertido a Python en una opción excelente para una gran variedad de aplicaciones. En muchos proyectos, el tiempo de desarrollo y mantenimiento resulta mucho más costoso que unos pocos milisegundos adicionales de ejecución.
En este volumen aprenderás los elementos fundamentales del lenguaje: cómo almacenar información, trabajar con diferentes tipos de datos, realizar operaciones y escribir programas sencillos utilizando condiciones, bucles y funciones.
\subsection{Uso}
Un programa en Python está formado por \textbf{instrucciones} que se ejecutan una tras otra. Cada instrucción indica al ordenador una acción concreta, como almacenar un valor, realizar un cálculo, mostrar un mensaje o repetir una tarea.
Python permite programar siguiendo distintos paradigmas, como la programación estructurada, la programación funcional o la programación orientada a objetos (POO). Esta flexibilidad permite elegir el enfoque más adecuado para cada problema y es una de las razones por las que Python resulta un lenguaje tan versátil.
A medida que avances en este libro descubrirás que, combinando unas pocas instrucciones básicas, es posible construir programas cada vez más completos y útiles.
No es necesario memorizar toda la sintaxis desde el principio. Lo importante es comprender qué hace cada elemento del lenguaje y cuándo utilizarlo.
\subsection{Ejemplo}
El siguiente programa muestra un mensaje por pantalla:


\begin{tcolorbox}
\scriptsize
 \begin{verbatim}
print("¡Hola, mundo!")
\end{verbatim}
\end{tcolorbox}
Aunque es un ejemplo muy sencillo, ya contiene varios conceptos importantes que iremos estudiando a lo largo del libro: una función ({\scriptsize \texttt{print}}), una cadena de texto ({\scriptsize \texttt{"¡Hola, mundo!"}}) y una instrucción que el intérprete ejecuta.
\subsection{Conclusión}
Python es un lenguaje sencillo, potente y versátil. Su sintaxis clara permite aprender a programar de forma progresiva, centrándose en la resolución de problemas en lugar de en la complejidad del lenguaje.
En los siguientes apartados aprenderás los bloques fundamentales del lenguaje y adquirirás las herramientas necesarias para comenzar a desarrollar tus propios programas en Python.

\newpage

\section{Variables}
\subsection{Introducción}
Los programas necesitan \textbf{recordar información} mientras se ejecutan. Por ejemplo, el nombre de un usuario, el resultado de un cálculo o el número de intentos realizados durante un juego.
Sin una forma de almacenar esa información, el programa tendría que volver a calcularla o solicitarla cada vez que la necesitara. Esto haría que los programas fueran poco prácticos e incluso imposibles de desarrollar.
Para resolver este problema existen las \textbf{variables}. Una variable permite asociar un nombre a un valor para poder utilizarlo más adelante dentro del programa.
\subsection{Uso}
En Python, una variable se crea mediante una \textbf{asignación}. Para ello, se escribe el nombre de la variable, seguido del operador de asignación ({\scriptsize \texttt{=}}) y el valor que se desea almacenar.


\begin{tcolorbox}
\scriptsize
 \begin{verbatim}
edad = 30
\end{verbatim}
\end{tcolorbox}
En una asignación, Python evalúa primero la expresión situada a la derecha del operador {\scriptsize \texttt{=}} y, una vez obtenido el resultado, lo asigna a la variable situada a la izquierda. Al asignar el valor de una variable a otra, la nueva variable pasa a contener el mismo valor.
Por ejemplo:


\begin{tcolorbox}
\scriptsize
 \begin{verbatim}
edad = 30
edad_pepe = edad
\end{verbatim}
\end{tcolorbox}
Hace que la variable edad\_pepe sea igual a la variable edad. Mientras que:


\begin{tcolorbox}
\scriptsize
 \begin{verbatim}
edad = 30
edad = edad_pepe # Error: edad_pepe todavía no existe.
\end{verbatim}
\end{tcolorbox}
Daría error porque no se ha definido previamente edad\_pepe y al intentar asignarle el valor de edad\_pepe a edad falta información.
A partir de ese momento, el nombre {\scriptsize \texttt{edad}} representa el valor almacenado y puede utilizarse en cualquier parte del programa.
El nombre de una variable debe ser descriptivo y representar con claridad la información que contiene. Elegir nombres adecuados facilita la lectura y el mantenimiento del código.


\begin{tcolorbox}
\scriptsize
 \begin{verbatim}
nombre = "Ana"
precio_total = 24.95
numero_intentos = 3
\end{verbatim}
\end{tcolorbox}
Además hay dos formas de asignación: directa e indirecta. La asignación directa es darle un valor concreto a la variable. Además de asignar un valor directamente, una variable también puede obtener su valor a partir de una expresión, del resultado de una función, de la entrada del usuario o del valor de otra variable. A esta forma de obtener el valor la llamaremos asignación indirecta.


\begin{tcolorbox}
\scriptsize
 \begin{verbatim}
precio = 12.5
cantidad = 4

total = precio * cantidad
\end{verbatim}
\end{tcolorbox}
\subsection{Ejemplo}
El siguiente programa almacena el nombre y la edad de una persona y posteriormente muestra esa información por pantalla.


\begin{tcolorbox}
\scriptsize
 \begin{verbatim}
nombre = "Ana"
edad = 30

print(nombre)
print(edad)
\end{verbatim}
\end{tcolorbox}
Las variables también pueden utilizarse para realizar operaciones.


\begin{tcolorbox}
\scriptsize
 \begin{verbatim}
precio = 12.5
cantidad = 4

total = precio * cantidad

print(total)
\end{verbatim}
\end{tcolorbox}
\subsection{Conclusión}
Las variables permiten almacenar información para reutilizarla durante la ejecución de un programa. Gracias a ellas, es posible recordar valores, realizar cálculos y construir programas que respondan a diferentes situaciones.
En el siguiente capítulo descubrirás que no toda la información es igual. Python puede almacenar números, textos, valores lógicos y otros muchos tipos de datos, cada uno diseñado para representar una clase diferente de información.

\newpage

\section{Tipos de datos}
\subsection{Introducción}
No toda la información que utiliza un programa es igual. Un número, un texto o una lista de elementos representan conceptos diferentes y permiten realizar operaciones distintas.
Por ejemplo, dos números pueden sumarse, mientras que un texto puede concatenarse con otro. Del mismo modo, una lista puede contener varios elementos, mientras que un número representa un único valor.
Para distinguir cada clase de información, Python utiliza los \textbf{tipos de datos}. Cada valor pertenece a un tipo determinado, que define cómo se almacena y qué operaciones pueden realizarse sobre él.
\subsection{Uso}
En Python, normalmente no es necesario indicar el tipo de una variable. Python determina automáticamente el tipo a partir del valor asignado.


\begin{tcolorbox}
\scriptsize
 \begin{verbatim}
nombre = "Ana"
edad = 30
altura = 1.68
estudiante = True
\end{verbatim}
\end{tcolorbox}
En este ejemplo, Python reconoce que {\scriptsize \texttt{30}} es un número entero, {\scriptsize \texttt{"Ana"}} es un texto y {\scriptsize \texttt{True}} representa un valor lógico.
Aunque todas las variables se crean mediante una asignación, el tipo del valor almacenado es diferente en cada una de ellas.
A lo largo de este libro estudiarás los tipos de datos más utilizados y aprenderás cuándo conviene emplear cada uno de ellos.
\subsection{Ejemplo}
El siguiente programa utiliza varios tipos de datos diferentes.


\begin{tcolorbox}
\scriptsize
 \begin{verbatim}
nombre = "Ana"
edad = 30
altura = 1.68
estudiante = True

print(nombre)
print(edad)
print(altura)
print(estudiante)
\end{verbatim}
\end{tcolorbox}
Aunque todas las variables se crean de la misma forma, cada una almacena un tipo de información diferente.
\subsection{Conclusión}
Los tipos de datos permiten representar correctamente la información que utiliza un programa. Cada tipo está diseñado para almacenar una clase concreta de valores y permite realizar un conjunto de operaciones adaptadas a ellos.
En los próximos capítulos estudiarás los principales tipos de datos de Python y aprenderás a utilizarlos en situaciones reales.

\newpage


\chapter{Tipos de datos}

\section{Números}
\subsection{Introducción}
Los números están presentes en la mayoría de los programas. Se utilizan para realizar cálculos, contar elementos, medir tiempos, representar precios o almacenar cualquier otra cantidad numérica.
Python dispone de distintos tipos de datos para representar números. Los más utilizados son los \textbf{enteros} ({\scriptsize \texttt{int}}) y los \textbf{decimales} ({\scriptsize \texttt{float}}).
Los enteros representan números sin parte decimal.


\begin{tcolorbox}
\scriptsize
 \begin{verbatim}
edad = 30
temperatura = -5
\end{verbatim}
\end{tcolorbox}
Los decimales representan números con parte fraccionaria.


\begin{tcolorbox}
\scriptsize
 \begin{verbatim}
altura = 1.75
precio = 12.95
\end{verbatim}
\end{tcolorbox}
\subsection{Uso}
Los números permiten realizar operaciones matemáticas de forma sencilla.


\begin{tcolorbox}
\scriptsize
 \begin{verbatim}
a = 10
b = 3

print(a + b)
print(a - b)
print(a * b)
print(a / b)
\end{verbatim}
\end{tcolorbox}
Además de las operaciones básicas, Python incorpora otros operadores muy útiles.


\begin{tcolorbox}
\scriptsize
 \begin{verbatim}
print(10 // 3)   # División entera
print(10 % 3)    # Resto
print(2 ** 3)    # Potencia
\end{verbatim}
\end{tcolorbox}
Las operaciones pueden combinarse para construir expresiones más complejas.


\begin{tcolorbox}
\scriptsize
 \begin{verbatim}
resultado = (5 + 3) * 2
\end{verbatim}
\end{tcolorbox}
Python respeta la precedencia habitual de los operadores matemáticos, aunque el uso de paréntesis mejora la claridad del código y evita errores.
Una característica interesante de Python es que los números enteros no tienen un tamaño máximo fijo. A diferencia de otros lenguajes, donde un entero suele ocupar un número determinado de bits, Python aumenta automáticamente el espacio necesario para almacenar valores cada vez mayores. En la práctica, el único límite es la memoria disponible del sistema.
\subsection{Ejemplo}
El siguiente ejemplo calcula el importe total de una compra.


\begin{tcolorbox}
\scriptsize
 \begin{verbatim}
precio = 19.95
cantidad = 3

total = precio * cantidad

print(total)
\end{verbatim}
\end{tcolorbox}
También es posible combinar varias operaciones en una misma expresión.


\begin{tcolorbox}
\scriptsize
 \begin{verbatim}
nota1 = 8
nota2 = 7
nota3 = 9

media = (nota1 + nota2 + nota3) / 3

print(media)
\end{verbatim}
\end{tcolorbox}
\subsection{Conclusión}
Los números permiten representar cantidades y realizar cálculos dentro de un programa. Python distingue principalmente entre números enteros ({\scriptsize \texttt{int}}) y números decimales ({\scriptsize \texttt{float}}), proporcionando operadores que permiten construir desde operaciones sencillas hasta expresiones matemáticas más complejas.
En el siguiente capítulo descubrirás otro tipo de dato fundamental: las cadenas de texto, utilizadas para representar información escrita.

\newpage

\section{Cadenas de texto}
\subsection{Introducción}
Los programas no solo trabajan con números. También necesitan representar información escrita, como nombres, direcciones, mensajes, contraseñas o el contenido de un archivo.
Para almacenar este tipo de información, Python utiliza las \textbf{cadenas de texto} ({\scriptsize \texttt{str}}). Una cadena es una secuencia de caracteres que puede contener letras, números, símbolos o espacios.
Las cadenas se delimitan mediante comillas simples ({\scriptsize \texttt{'}}) o dobles ({\scriptsize \texttt{"}}), siendo ambas equivalentes en Python.


\begin{tcolorbox}
\scriptsize
 \begin{verbatim}
nombre = "Ana"
mensaje = '¡Hola, mundo!'
\end{verbatim}
\end{tcolorbox}
\subsection{Uso}
Las cadenas pueden almacenarse en variables y utilizarse como cualquier otro tipo de dato.


\begin{tcolorbox}
\scriptsize
 \begin{verbatim}
ciudad = "Madrid"
\end{verbatim}
\end{tcolorbox}
También es posible combinar varias cadenas mediante el operador de concatenación ({\scriptsize \texttt{+}}).


\begin{tcolorbox}
\scriptsize
 \begin{verbatim}
nombre = "Ana"
apellido = "García"

nombre_completo = nombre + " " + apellido
\end{verbatim}
\end{tcolorbox}
El operador de multiplicación ({\scriptsize \texttt{*}}) permite repetir una cadena un número determinado de veces.


\begin{tcolorbox}
\scriptsize
 \begin{verbatim}
separador = "-" * 20
\end{verbatim}
\end{tcolorbox}
Python distingue entre mayúsculas y minúsculas, por lo que dos cadenas que solo se diferencian en el uso de estas letras se consideran distintas.


\begin{tcolorbox}
\scriptsize
 \begin{verbatim}
usuario = "ana"
administrador = "Ana"
\end{verbatim}
\end{tcolorbox}
\subsection{Ejemplo}
El siguiente programa crea un mensaje de bienvenida utilizando varias cadenas de texto.


\begin{tcolorbox}
\scriptsize
 \begin{verbatim}
nombre = "Ana"

mensaje = "Bienvenida, " + nombre + "."

print(mensaje)
\end{verbatim}
\end{tcolorbox}
También es posible repetir una cadena para generar un separador.


\begin{tcolorbox}
\scriptsize
 \begin{verbatim}
print("=" * 30)
print("Fin del programa")
print("=" * 30)
\end{verbatim}
\end{tcolorbox}
\subsection{Conclusión}
Las cadenas de texto permiten representar información escrita dentro de un programa. Gracias a ellas es posible almacenar, mostrar y combinar texto de forma sencilla.
En el siguiente capítulo conocerás los \textbf{valores lógicos} ({\scriptsize \texttt{bool}}), un tipo de dato utilizado para representar situaciones que solo pueden tener dos posibles estados: verdadero o falso.

\newpage

\section{Booleanos}
\subsection{Introducción}
En muchas situaciones, un programa necesita responder a preguntas cuya respuesta solo puede tener dos posibilidades: \textbf{sí} o \textbf{no}.
Por ejemplo:

\begin{itemize}
\item  ¿El usuario ha iniciado sesión?
\item  ¿La contraseña es correcta?
\item  ¿El pedido ha sido enviado?
\item  ¿El número es mayor que cero?
\end{itemize}
Para representar este tipo de información, Python utiliza el tipo de dato \textbf{booleano} ({\scriptsize \texttt{bool}}), cuyos únicos valores posibles son {\scriptsize \texttt{True}} (verdadero) y {\scriptsize \texttt{False}} (falso).
\subsection{Uso}
Los valores booleanos pueden asignarse a una variable como cualquier otro tipo de dato.


\begin{tcolorbox}
\scriptsize
 \begin{verbatim}
es_mayor_de_edad = True
sesion_iniciada = False
\end{verbatim}
\end{tcolorbox}
También pueden obtenerse como resultado de una comparación.


\begin{tcolorbox}
\scriptsize
 \begin{verbatim}
edad = 20

es_adulto = edad >= 18
\end{verbatim}
\end{tcolorbox}
En este ejemplo, la comparación produce el valor {\scriptsize \texttt{True}}, que queda almacenado en la variable {\scriptsize \texttt{es\_adulto}}.
Los valores booleanos suelen utilizarse para controlar el comportamiento de un programa, permitiendo tomar decisiones o repetir acciones, como veremos en capítulos posteriores.
\subsection{Ejemplo}
El siguiente programa comprueba si una persona es mayor de edad.


\begin{tcolorbox}
\scriptsize
 \begin{verbatim}
edad = 20

es_adulto = edad >= 18

print(es_adulto)
\end{verbatim}
\end{tcolorbox}
También es posible asignar directamente un valor booleano.


\begin{tcolorbox}
\scriptsize
 \begin{verbatim}
modo_oscuro = True

print(modo_oscuro)
\end{verbatim}
\end{tcolorbox}
\subsection{Conclusión}
Los booleanos permiten representar información que solo puede adoptar dos estados: verdadero o falso. Son fundamentales para expresar condiciones y controlar el comportamiento de un programa.
En el siguiente capítulo descubrirás las \textbf{listas}, un tipo de dato que permite almacenar varios valores en una única variable.

\newpage

\section{Listas}
\subsection{Introducción}
Hasta ahora hemos trabajado con variables capaces de almacenar un único valor. Sin embargo, en muchas ocasiones es necesario trabajar con varios valores relacionados.
Por ejemplo, un programa puede necesitar almacenar las notas de un alumno, los nombres de los participantes de un curso o los productos de un carrito de compra.
Crear una variable para cada dato sería poco práctico. Para resolver este problema, Python dispone de las \textbf{listas} ({\scriptsize \texttt{list}}), un tipo de dato que permite almacenar varios valores en una única variable.
Las listas mantienen el orden de los elementos, pueden contener valores del mismo tipo o de tipos diferentes y, además, son modificables.
\subsection{Uso}
Las listas se representan mediante corchetes ({\scriptsize \texttt{[]}}), separando cada elemento mediante comas.


\begin{tcolorbox}
\scriptsize
 \begin{verbatim}
numeros = [3, 7, 12, 25]
\end{verbatim}
\end{tcolorbox}
También pueden almacenar cadenas de texto.


\begin{tcolorbox}
\scriptsize
 \begin{verbatim}
nombres = ["Ana", "Luis", "Marta"]
\end{verbatim}
\end{tcolorbox}
Incluso es posible combinar distintos tipos de datos en una misma lista.


\begin{tcolorbox}
\scriptsize
 \begin{verbatim}
datos = ["Ana", 30, True]
\end{verbatim}
\end{tcolorbox}
Cada elemento ocupa una posición dentro de la lista, denominada \textbf{índice}. En Python, los índices comienzan en {\scriptsize \texttt{0}}.


\begin{tcolorbox}
\scriptsize
 \begin{verbatim}
colores = ["rojo", "verde", "azul"]

print(colores[0])
print(colores[1])
print(colores[2])
\end{verbatim}
\end{tcolorbox}
El número de elementos de una lista puede variar durante la ejecución del programa.
\subsection{Ejemplo}
El siguiente programa almacena las temperaturas registradas durante una semana y muestra la primera y la última.


\begin{tcolorbox}
\scriptsize
 \begin{verbatim}
temperaturas = [18.5, 20.1, 19.8, 21.3, 22.0, 20.7, 19.4]

print(temperaturas[0])
print(temperaturas[6])
\end{verbatim}
\end{tcolorbox}
Las listas también permiten modificar un elemento indicando su índice.


\begin{tcolorbox}
\scriptsize
 \begin{verbatim}
temperaturas = [18.5, 20.1, 19.8]

temperaturas[1] = 21.0

print(temperaturas)
\end{verbatim}
\end{tcolorbox}
\subsection{Conclusión}
Las listas permiten agrupar varios valores relacionados en una única variable, manteniendo el orden en el que fueron almacenados. Gracias a ellas, es posible trabajar de forma cómoda con colecciones de datos.
En el siguiente capítulo descubrirás las \textbf{tuplas}, una estructura muy similar a las listas, pero diseñada para almacenar datos que no deben modificarse.

\newpage

\section{Tuplas}
\subsection{Introducción}
Las listas permiten almacenar varios valores relacionados y modificar su contenido cuando es necesario. Sin embargo, hay situaciones en las que los datos no deben cambiar una vez creados.
Por ejemplo, las coordenadas de un punto, los meses del año o los días de la semana son datos que permanecen constantes durante la ejecución de un programa.
Para representar este tipo de información, Python dispone de las \textbf{tuplas} ({\scriptsize \texttt{tuple}}), una colección ordenada cuyos elementos no pueden modificarse una vez creada.
\subsection{Uso}
Las tuplas se representan mediante paréntesis ({\scriptsize \texttt{()}}), separando sus elementos mediante comas.


\begin{tcolorbox}
\scriptsize
 \begin{verbatim}
coordenadas = (10, 20)
\end{verbatim}
\end{tcolorbox}
Al igual que las listas, pueden contener valores del mismo tipo o de tipos diferentes.


\begin{tcolorbox}
\scriptsize
 \begin{verbatim}
persona = ("Ana", 30, True)
\end{verbatim}
\end{tcolorbox}
Los elementos de una tupla también se acceden mediante índices, comenzando por {\scriptsize \texttt{0}}.


\begin{tcolorbox}
\scriptsize
 \begin{verbatim}
dias = ("Lunes", "Martes", "Miércoles")

print(dias[0])
print(dias[1])
print(dias[2])
\end{verbatim}
\end{tcolorbox}
Sin embargo, una vez creada la tupla, sus elementos no pueden modificarse.


\begin{tcolorbox}
\scriptsize
 \begin{verbatim}
dias = ("Lunes", "Martes", "Miércoles")

dias[0] = "Domingo"   # Error
\end{verbatim}
\end{tcolorbox}
\subsection{Ejemplo}
El siguiente programa almacena las coordenadas de un punto utilizando una tupla.


\begin{tcolorbox}
\scriptsize
 \begin{verbatim}
punto = (12, 8)

print(punto[0])
print(punto[1])
\end{verbatim}
\end{tcolorbox}
Las tuplas son especialmente útiles cuando un conjunto de datos representa una única entidad y no debe modificarse durante la ejecución del programa.
\subsection{Conclusión}
Las tuplas permiten agrupar varios valores relacionados de forma ordenada, igual que las listas. Su principal diferencia es que su contenido no puede modificarse una vez creado, lo que ayuda a proteger datos que deben permanecer constantes.
Elegir entre una lista y una tupla depende de si los datos deben cambiar durante la ejecución del programa.
En el siguiente capítulo conocerás los \textbf{conjuntos}, una colección diseñada para almacenar elementos sin repetirlos.

\newpage

\section{Conjuntos}
\subsection{Introducción}
Las listas y las tuplas permiten almacenar colecciones de elementos manteniendo el orden en el que fueron añadidos. Sin embargo, en algunas situaciones el orden no es importante, mientras que evitar elementos repetidos sí lo es.
Por ejemplo, un programa puede necesitar almacenar los idiomas que habla una persona, las etiquetas de un artículo o los permisos asignados a un usuario. En estos casos, un mismo elemento solo debería aparecer una vez.
Para resolver este problema, Python dispone de los \textbf{conjuntos} ({\scriptsize \texttt{set}}), una colección de elementos únicos en la que el orden de los elementos no está garantizado.
\subsection{Uso}
Los conjuntos se representan mediante llaves ({\scriptsize \texttt{{}}}), separando sus elementos mediante comas.
Un conjunto vacío no se crea con {\scriptsize \texttt{{}}}, ya que esa sintaxis corresponde a un diccionario vacío. Para crear un conjunto vacío se utiliza {\scriptsize \texttt{set()}}.


\begin{tcolorbox}
\scriptsize
 \begin{verbatim}
conjunto = set()

idiomas = {"Español", "Inglés", "Francés"}
\end{verbatim}
\end{tcolorbox}
Si un mismo elemento aparece varias veces, el conjunto solo conservará una copia.


\begin{tcolorbox}
\scriptsize
 \begin{verbatim}
numeros = {1, 2, 2, 3, 3, 3}

print(numeros)
\end{verbatim}
\end{tcolorbox}
La salida será:


\begin{tcolorbox}
\scriptsize
 \begin{verbatim}
{1, 2, 3}
\end{verbatim}
\end{tcolorbox}
Al no mantener un orden definido, los conjuntos no permiten acceder a sus elementos mediante índices.


\begin{tcolorbox}
\scriptsize
 \begin{verbatim}
idiomas = {"Español", "Inglés", "Francés"}

print(idiomas[0])   # Error
\end{verbatim}
\end{tcolorbox}
\subsection{Ejemplo}
El siguiente programa elimina automáticamente los nombres repetidos de una colección.


\begin{tcolorbox}
\scriptsize
 \begin{verbatim}
participantes = {
    "Ana",
    "Luis",
    "Ana",
    "María",
    "Luis"
}

print(participantes)
\end{verbatim}
\end{tcolorbox}
El conjunto contendrá únicamente los nombres diferentes.
\subsection{Conclusión}
Los conjuntos permiten almacenar colecciones de elementos sin duplicados. Son especialmente útiles cuando lo importante es saber qué elementos existen, cuando lo importante es la presencia de un elemento y no la posición que ocupa.
En el siguiente capítulo descubrirás los \textbf{diccionarios}, una colección que permite asociar cada valor a una clave para acceder a la información de una forma más sencilla.

\newpage

\section{Diccionarios}
\subsection{Introducción}
Hasta ahora hemos trabajado con colecciones cuyos elementos se identifican por su posición o simplemente por su pertenencia al conjunto. Sin embargo, en muchas ocasiones resulta más útil asociar cada dato a un nombre que describa su significado.
Por ejemplo, al representar una persona es más natural hablar de su nombre, edad o ciudad que del primer, segundo o tercer elemento de una colección.
Para resolver este problema, Python dispone de los \textbf{diccionarios} ({\scriptsize \texttt{dict}}), una colección que almacena pares \textbf{clave-valor} (\textit{key-value pairs}), permitiendo acceder a cada dato mediante una clave en lugar de un índice.
\subsection{Uso}
Los diccionarios se representan mediante llaves ({\scriptsize \texttt{{}}}). Cada par clave-valor está formado por una clave y un valor, separados por dos puntos ({\scriptsize \texttt{:}}).


\begin{tcolorbox}
\scriptsize
 \begin{verbatim}
persona = {
    "nombre": "Ana",
    "edad": 30,
    "ciudad": "Madrid"
}
\end{verbatim}
\end{tcolorbox}
Para acceder a un valor se utiliza su clave.


\begin{tcolorbox}
\scriptsize
 \begin{verbatim}
print(persona["nombre"])
print(persona["edad"])
\end{verbatim}
\end{tcolorbox}
También es posible modificar el valor asociado a una clave existente.


\begin{tcolorbox}
\scriptsize
 \begin{verbatim}
persona["edad"] = 31
\end{verbatim}
\end{tcolorbox}
O añadir un nuevo par clave-valor.


\begin{tcolorbox}
\scriptsize
 \begin{verbatim}
persona["profesion"] = "Ingeniera"
\end{verbatim}
\end{tcolorbox}
Las claves de un diccionario deben ser únicas. Si una misma clave aparece varias veces, solo se conservará la última asignación.
Además de acceder a un valor mediante su clave, también es posible trabajar con el conjunto de claves o con el conjunto de valores de un diccionario. Más adelante aprenderás a obtener las claves mediante {\scriptsize \texttt{keys()}} y los valores mediante {\scriptsize \texttt{values()}}.
\subsection{Ejemplo}
El siguiente programa almacena información sobre un libro y muestra algunos de sus datos.


\begin{tcolorbox}
\scriptsize
 \begin{verbatim}
libro = {
    "titulo": "Python básico",
    "autor": "Ana García",
    "paginas": 250
}

print(libro["titulo"])
print(libro["paginas"])
\end{verbatim}
\end{tcolorbox}
Modificar un valor resulta tan sencillo como realizar una nueva asignación sobre su clave.


\begin{tcolorbox}
\scriptsize
 \begin{verbatim}
libro["paginas"] = 275

print(libro)
\end{verbatim}
\end{tcolorbox}
\subsection{Conclusión}
Los diccionarios permiten organizar la información mediante pares clave-valor, facilitando el acceso a cada dato por su nombre en lugar de por su posición. Son especialmente útiles para representar entidades con diferentes características o atributos.
Con este capítulo concluye el estudio de los principales tipos de datos de Python. En los siguientes capítulos aprenderás a utilizarlos para realizar operaciones, tomar decisiones, repetir tareas y organizar mejor tus programas.

\newpage


\chapter{Operadores y conversiones}

\section{Operadores aritméticos}
\subsection{Introducción}
Los números permiten representar cantidades, pero su utilidad reside en poder realizar operaciones con ellos.
Python incorpora distintos \textbf{operadores aritméticos} que permiten efectuar cálculos de forma sencilla, desde sumas y restas hasta divisiones, potencias y otras operaciones matemáticas.
\subsection{Uso}
Los operadores aritméticos más utilizados son:

\newpage

\begin{tcolorbox}
\begin{tabular}{ll}
 Operador  &  Operación        \\
 {\scriptsize \texttt{+}}       &  Suma             \\
 {\scriptsize \texttt{-}}       &  Resta            \\
 {\scriptsize \texttt{*}}       &  Multiplicación   \\
 {\scriptsize \texttt{/}}       &  División         \\
 {\scriptsize \texttt{//}}      &  División entera  \\
 {\scriptsize \texttt{\%}}       &  Resto (módulo)   \\
 {\scriptsize \texttt{**}}      &  Potencia         \\
\end{tabular}
\end{tcolorbox}
Las operaciones pueden combinarse para formar expresiones más complejas.


\begin{tcolorbox}
\scriptsize
 \begin{verbatim}
resultado = (8 + 2) * 5
\end{verbatim}
\end{tcolorbox}
Python sigue las reglas habituales de precedencia matemática, aunque el uso de paréntesis mejora la claridad del código.
\subsection{Ejemplo}
El siguiente programa calcula las horas, los minutos y los segundos a partir del número total de segundos transcurridos.

\newpage


\begin{tcolorbox}
\scriptsize
 \begin{verbatim}
segundos = 3675

horas = segundos // 3600
segundos_restantes = segundos % 3600

minutos = segundos_restantes // 60
segundos_finales  = segundos_restantes % 60

print("Horas:", horas)
print("Minutos:", minutos)
print("Segundos:", segundos_finales )
\end{verbatim}
\end{tcolorbox}
La salida será:


\begin{tcolorbox}
\scriptsize
 \begin{verbatim}
Horas: 1
Minutos: 1
Segundos: 15
\end{verbatim}
\end{tcolorbox}
En este ejemplo, la \textbf{división entera} ({\scriptsize \texttt{//}}) permite obtener el número de horas y minutos completos, mientras que el \textbf{operador resto o módulo} ({\scriptsize \texttt{\%}}) obtiene los segundos que aún quedan por convertir.
\subsection{Conclusión}
Los operadores aritméticos permiten realizar cálculos utilizando valores numéricos y construir expresiones cada vez más complejas.
En el siguiente capítulo aprenderás los \textbf{operadores de comparación}, que permiten comprobar relaciones entre valores y obtener un resultado lógico.

\newpage

\section{Operadores de comparación}
\subsection{Introducción}
En muchas ocasiones no basta con almacenar información o realizar cálculos. Un programa también necesita \textbf{comparar valores} para saber si una determinada condición se cumple.
Por ejemplo:

\begin{itemize}
\item  ¿Es una persona mayor de edad?
\item  ¿La contraseña introducida coincide con la esperada?
\item  ¿Dos números son iguales?
\item  ¿Un producto cuesta más que otro?
\end{itemize}
Para responder a estas preguntas, Python dispone de los \textbf{operadores de comparación}, que comparan dos valores y producen un resultado booleano: {\scriptsize \texttt{True}} o {\scriptsize \texttt{False}}.
\subsection{Uso}
Los operadores de comparación más utilizados son:

\newpage

\begin{tcolorbox}
\begin{tabular}{ll}
 Operador  &  Comparación        \\
 {\scriptsize \texttt{==}}      &  Igual que          \\
 {\scriptsize \texttt{!=}}      &  Distinto de        \\
 {\scriptsize \texttt{\textgreater{}}}       &  Mayor que          \\
 {\scriptsize \texttt{\textless{}}}       &  Menor que          \\
 {\scriptsize \texttt{\textgreater{}=}}      &  Mayor o igual que  \\
 {\scriptsize \texttt{\textless{}=}}      &  Menor o igual que  \\
\end{tabular}
\end{tcolorbox}
El resultado de una comparación siempre es un valor booleano.


\begin{tcolorbox}
\scriptsize
 \begin{verbatim}
edad = 20

print(edad >= 18)
\end{verbatim}
\end{tcolorbox}
También es posible almacenar el resultado de una comparación en una variable.


\begin{tcolorbox}
\scriptsize
 \begin{verbatim}
precio = 12.50

es_caro = precio > 10

print(es_caro)
\end{verbatim}
\end{tcolorbox}
Es importante no confundir el operador de asignación ({\scriptsize \texttt{=}}), utilizado para almacenar un valor en una variable, con el operador de comparación ({\scriptsize \texttt{==}}), que comprueba si dos valores son iguales.


\begin{tcolorbox}
\scriptsize
 \begin{verbatim}
edad = 18      # Asigna el valor 18 a la variable
edad == 18     # Comprueba si edad es igual a 18
\end{verbatim}
\end{tcolorbox}
\subsection{Ejemplo}
El siguiente programa comprueba si un alumno ha aprobado un examen.


\begin{tcolorbox}
\scriptsize
 \begin{verbatim}
nota = 7.5

aprobado = nota >= 5

print(aprobado)
\end{verbatim}
\end{tcolorbox}
También es posible comparar cadenas de texto.


\begin{tcolorbox}
\scriptsize
 \begin{verbatim}
usuario = "Ana"

print(usuario == "Ana")
\end{verbatim}
\end{tcolorbox}
\subsection{Conclusión}
Los operadores de comparación permiten comprobar la relación entre dos valores y obtener un resultado booleano. Son la base para que un programa pueda tomar decisiones en función de la información que procesa.
En el siguiente capítulo aprenderás los \textbf{operadores lógicos}, que permiten combinar varias condiciones para construir expresiones más complejas.

\newpage

\section{Operadores lógicos}
\subsection{Introducción}
En ocasiones, una única comparación no es suficiente para tomar una decisión. Un programa puede necesitar comprobar que se cumplen varias condiciones al mismo tiempo o que basta con que se cumpla una de ellas.
Por ejemplo:
* ¿Es mayor de edad \textbf{y} tiene entrada?
* ¿Es administrador \textbf{o} moderador?
* ¿El usuario \textbf{no} ha iniciado sesión?
Para resolver estas situaciones, Python dispone de los \textbf{operadores lógicos}, que permiten combinar o modificar valores booleanos.
\subsection{Uso}
Los operadores lógicos más utilizados son:

\newpage

\begin{tcolorbox}
\begin{tabular}{ll}
 \textbf{Operador}  &  \textbf{Significado}                                          \\
 {\scriptsize \texttt{and}}     &  Todas las condiciones deben ser verdaderas.          \\
 {\scriptsize \texttt{or}}      &  Basta con una de las condiciones verdadera.          \\
 {\scriptsize \texttt{not}}     &  Invierte el valor lógico de una condición.           \\
\end{tabular}
\end{tcolorbox}
Los operadores lógicos trabajan con valores booleanos ({\scriptsize \texttt{True}} y {\scriptsize \texttt{False}}) y producen también un resultado booleano.


\begin{tcolorbox}
\scriptsize
 \begin{verbatim}
es_adulto = True
tiene_entrada = False

print(es_adulto and tiene_entrada)
print(es_adulto or tiene_entrada)
print(not tiene_entrada)
\end{verbatim}
\end{tcolorbox}
\subsection{Ejemplo}
El siguiente programa comprueba distintas condiciones para permitir el acceso a un evento.


\begin{tcolorbox}
\scriptsize
 \begin{verbatim}
edad = 20
tiene_entrada = True
esta_en_lista = False

puede_entrar = edad >= 18 and tiene_entrada

print(puede_entrar)

acceso_especial = tiene_entrada or esta_en_lista

print(acceso_especial)

print(not esta_en_lista)
\end{verbatim}
\end{tcolorbox}
En este ejemplo:

\begin{itemize}
\item  {\scriptsize \texttt{and}} exige que ambas condiciones sean verdaderas.
\item  {\scriptsize \texttt{or}} devuelve {\scriptsize \texttt{True}} si al menos una condición es verdadera.
\item  {\scriptsize \texttt{not}} invierte el valor lógico de una condición.
\end{itemize}
\subsection{Conclusión}
Los operadores lógicos permiten combinar o modificar condiciones para construir expresiones más complejas. Gracias a ellos, un programa puede comprobar simultáneamente varios requisitos o invertir el resultado de una condición.
En el siguiente capítulo descubrirás cómo convertir valores entre los distintos tipos de datos de Python, una herramienta fundamental para trabajar con información de diferente naturaleza.

\newpage

\section{Conversión de tipos}
\subsection{Introducción}
Hasta ahora hemos trabajado con distintos tipos de datos, como números, cadenas de texto, valores booleanos o colecciones. En muchas ocasiones, un programa necesita transformar un valor de un tipo a otro para poder seguir trabajando con él.
Por ejemplo, un número puede convertirse en una cadena de texto para mostrarlo por pantalla, o un texto que contiene un número puede convertirse en un valor numérico para realizar cálculos.
Para realizar estas transformaciones, Python proporciona varias funciones de conversión de tipos.
\subsection{Uso}
Las conversiones más habituales se realizan mediante las funciones {\scriptsize \texttt{int()}}, {\scriptsize \texttt{float()}}, {\scriptsize \texttt{str()}} y {\scriptsize \texttt{bool()}}.


\begin{tcolorbox}
\scriptsize
 \begin{verbatim}
numero = int("25")
precio = float("19.95")
texto = str(42)
activo = bool(1)
\end{verbatim}
\end{tcolorbox}
Cada una de estas funciones crea un nuevo valor del tipo indicado a partir del valor original.
No todas las conversiones son posibles. Por ejemplo, intentar convertir un texto que no representa un número producirá un error.


\begin{tcolorbox}
\scriptsize
 \begin{verbatim}
numero = int("Hola")    # Error
\end{verbatim}
\end{tcolorbox}
Por este motivo, es importante asegurarse de que el contenido puede convertirse al tipo deseado.
\subsection{Ejemplo}
El siguiente programa convierte la edad almacenada como una cadena de texto en un número entero para poder realizar una operación.


\begin{tcolorbox}
\scriptsize
 \begin{verbatim}
edad_texto = "30"

edad = int(edad_texto)

print(edad + 5)
\end{verbatim}
\end{tcolorbox}
También es posible convertir un número en una cadena de texto.


\begin{tcolorbox}
\scriptsize
 \begin{verbatim}
temperatura = 21.5

mensaje = "La temperatura es " + str(temperatura) + " grados."

print(mensaje)
\end{verbatim}
\end{tcolorbox}
\subsection{Conclusión}
La conversión de tipos permite adaptar la información al formato que necesita un programa en cada momento. Gracias a estas funciones es posible transformar números, textos y valores booleanos para realizar cálculos, mostrar información o combinar datos de diferentes tipos.
Con este capítulo concluye el estudio de los fundamentos de Python. A lo largo de este volumen has aprendido cómo almacenar información, representarla mediante distintos tipos de datos, organizarla en colecciones, realizar operaciones y convertirla entre diferentes formatos.
En el siguiente volumen utilizarás todas estas herramientas para construir programas capaces de tomar decisiones, repetir tareas y organizar el código mediante funciones.

\newpage


\chapter{Programación Orientada a Objetos}

\section{Programación Orientada a Objetos}
\subsection{Introducción}
En la programación clásica, el paradigma que más se repetía era la programación estructurada. Consistía en escribir código con instrucciones que se ejecutaban una detrás de otra. La Programación Orientada a Objetos (POO) consiste en el uso de Objetos como paradigma central. 
La programación orientada a objetos surgió como una forma de modelar sistemas complejos mediante objetos que representan entidades del problema. En lugar de organizar el programa únicamente como una secuencia de instrucciones, organiza el código alrededor de objetos que encapsulan datos y comportamiento.
Para ello utiliza clases, que definen la estructura y el comportamiento de un tipo de objeto mediante atributos y métodos. Un objeto es una instancia concreta de esa clase cuyos atributos poseen valores determinados.
\subsection{Uso de la POO}
La POO requiere de la creación de clases. Estas clases contienen variables llamadas atributos que al instanciar los objetos obtienen un valor propio de ese objeto y unos métodos que son funciones que hacen uso de los atributos internos del objeto junto con parámetros que se le pasen.
Una vez se tiene programada la clase, los objetos pueden modificar los atributos y emplear los métodos de la clase para poder realizar acciones, modificaciones o almacenar información.
La POO se compone de varios tipos de elementos:

\begin{itemize}
\item  las clases
\item  los atributos
\item  los métodos
\end{itemize}
Las clases son el modelo base de los objetos. Los atributos son variables que almacenan el estado de un objeto. Y los métodos son funciones propias de la clase que usan los objetos empleando atributos o no.
En cuanto a las bases de la POO, tiene cuatro pilares:

\begin{itemize}
\item  La abstracción
\item  La herencia
\item  La encapsulación
\item  El polimirfismo
\end{itemize}
La herencia es un proceso de especialización de las clases. La abstracción consiste en representar una entidad del problema mediante una clase que recoge únicamente los atributos y métodos relevantes para el contexto. La encapsulación consiste en controlar el acceso al estado interno de un objeto mediante una interfaz pública. Finalmente, el polimorfismo es el poder utilizar una misma interfaz para trabajar con objetos de distintos tipos.
Además de estos pilares, la composición es un mecanismo fundamental muy utilizado para construir objetos complejos reutilizando otros objetos.
\subsection{Ejemplo}
A lo largo de esta colección se utilizarán fragmentos reales del código de Odix. En este capítulo no es necesario comprender todos los detalles; su propósito es mostrar el aspecto que tiene una clase real. Cada uno de los conceptos empleados se estudiará en los capítulos siguientes.

\newpage


\begin{tcolorbox}
\scriptsize
 \begin{verbatim}
class Node(ABC):
    """Base class for all Tabula nodes.

    Every element of the AST inherits from this class. A node stores its
    position in the tree, its children and its unique identifier.

    Attributes:
        id: Unique node identifier.
        parent: Parent node. ``None`` if this node is the root.
        children: Child nodes.
    """

    def __init__(self, node_id: str) -> None:
        """Initializes a new node.

        Args:
            node_id: Unique identifier of the node.
        """
        self.id = node_id
        self.parent: Node | None = None
        self.children: list[Node] = []

    def add_child(self, child: Node) -> None:
        """Adds a child node.

        Args:
            child: Node to attach.
        """
        child.parent = self
        self.children.append(child)

    def remove_child(self, child: Node) -> None:
        """Removes a child node.

        Args:
            child: Child node to remove.

        Raises:
            ValueError: If the node is not a child of this node.
        """
        self.children.remove(child)
        child.parent = None
\end{verbatim}
\end{tcolorbox}
En esta clase aparecen varios de los conceptos de POO y de otros temas. Al final de la colección el lector podrá escribir su propio código con todos los conceptos incluidos en el código de Odix.
\subsection{Conclusión}
La programación orientada a objetos permite organizar el software mediante clases y objetos que encapsulan datos y comportamiento. A partir de estos conceptos se construyen programas más modulares, reutilizables y fáciles de mantener. En los siguientes capítulos se estudiará cada uno de los elementos fundamentales que forman este paradigma.

\newpage

\section{clases}
\subsection{Introducción}
Las clases son un tipo de elemento versátil usado en el paradigma de programación orientada a objetos (POO por sus siglas, que en inglés serían OOP). Python es un lenguaje multiparadigma. Aunque es posible programar utilizando únicamente funciones, también ofrece un potente modelo de programación orientada a objetos basado en clases.
Una clase es resultado del proceso de abstracción, en el que de uno o varios elementos se extraen los atributos y métodos principales y se agrupan sus atributos y comportamientos en un nuevo tipo de objeto. 
Un ejemplo de tipo integrado es {\scriptsize \texttt{int}}. Aunque habitualmente se piense en él simplemente como el tipo de los números enteros, en realidad {\scriptsize \texttt{int}} es una clase. Cada número entero ({\scriptsize \texttt{1}}, {\scriptsize \texttt{2}}, {\scriptsize \texttt{3}}, ...) es una instancia de dicha clase y dispone de operaciones propias, como la suma, la resta o la multiplicación, implementadas mediante los métodos especiales del tipo.
Como se ha explicado, un objeto es la instancia de una clase. Se pueden crear tantos objetos como se necesiten de una clase para llevar a cabo la ejecución del programa. Cada objeto posee sus propios atributos, mientras que comparte los métodos definidos por la clase. El objeto no añade características a la clase sino que la utiliza como patrón para ser creado y es único.
Las clases se emplean para simplificar y agrupar elementos que tienen un comportamiento similar con operaciones propias y diferentes de las que vienen por defecto para variables existentes. Suelen ser elementos más complejos que las variables básicas y suelen contener internamente como atributos variables u otras clases.
\subsection{Uso de las clases}
Para definir una clase en Python se emplea la palabra clave {\scriptsize \texttt{class}}, seguida del nombre de la clase y dos puntos ({\scriptsize \texttt{:}}). Como cualquier otro bloque del lenguaje, el contenido de la clase se escribe con un nivel de indentación. Aunque el nombre puede ser el que el programador desee, se recomienda que sea descriptivo y siga la convención \textbf{PascalCase}, es decir, comenzando cada palabra por una letra mayúscula. Esta convención permite identificar rápidamente que se trata de una clase y diferenciarla de las funciones y variables, que normalmente se escriben en minúsculas utilizando \textit{snake\_case}.
Dentro de la clase suele definirse, en primer lugar, el constructor. El constructor es un método especial denominado {\scriptsize \texttt{\_\_init\_\_}} que Python ejecuta automáticamente cuando se crea un nuevo objeto. Su función consiste en inicializar el estado del objeto, asignando los valores iniciales de sus atributos a partir de los parámetros recibidos.
Los atributos representan la información que almacena cada objeto. Habitualmente se crean dentro del constructor mediante la sintaxis {\scriptsize \texttt{self.atributo}}, aunque también pueden existir atributos de clase, que son compartidos por todas las instancias y se definen directamente dentro de la clase, fuera de cualquier método.
Además de almacenar información, las clases definen métodos. Los métodos son funciones asociadas a la clase que describen el comportamiento de sus objetos. Al invocarlos, pueden modificar el estado del objeto, consultar información o realizar cualquier otra operación relacionada con él.
Todos los métodos de instancia reciben como primer parámetro una referencia al propio objeto. Por convención, este parámetro recibe el nombre {\scriptsize \texttt{self}}. Aunque podría utilizarse cualquier identificador, el uso de self está completamente estandarizado en Python y debe respetarse. Gracias a este parámetro, un método puede acceder a los atributos y a otros métodos del mismo objeto sin necesidad de recibirlos como argumentos adicionales. En otras palabras, {\scriptsize \texttt{self}} representa el objeto sobre el que se está ejecutando el método.
Una vez definida la clase, es posible crear tantos objetos como sean necesarios. Este proceso recibe el nombre de instanciación y consiste en llamar a la clase como si fuera una función, proporcionando los argumentos que requiera su constructor.
Para acceder a los atributos o invocar un método de un objeto se utiliza el operador punto ({\scriptsize \texttt{.}}). Los atributos se consultan escribiendo {\scriptsize \texttt{objeto.atributo}}, mientras que los métodos se ejecutan mediante {\scriptsize \texttt{objeto.metodo()}}, proporcionando entre paréntesis los argumentos que dicho método necesite.
Por último, una clase puede heredar atributos y métodos de otra clase. Este mecanismo, conocido como \textbf{herencia}, permite reutilizar código y especializar comportamientos. En Python, la herencia se indica escribiendo el nombre de la clase padre entre paréntesis tras el nombre de la clase hija, por ejemplo: {\scriptsize \texttt{class claseHija(clasePadre):}}.
\subsection{Ejemplo}
Esta es la clase Node del proyecto que inspira la colección de libros de Python Real: 

\newpage


\begin{tcolorbox}
\scriptsize
 \begin{verbatim}
class Node(ABC):
    """Base class for all Tabula nodes.

    Every element of the AST inherits from this class. A node stores its
    position in the tree, its children and its unique identifier.

    Attributes:
        id: Unique node identifier.
        parent: Parent node. ``None`` if this node is the root.
        children: Child nodes.
    """

    def __init__(self, node_id: str) -> None:
        """Initializes a new node.

        Args:
            node_id: Unique identifier of the node.
        """
        self.id = node_id
        self.parent: Node | None = None
        self.children: list[Node] = []

    def add_child(self, child: Node) -> None:
        """Adds a child node.

        Args:
            child: Node to attach.
        """
        child.parent = self
        self.children.append(child)

    def remove_child(self, child: Node) -> None:
        """Removes a child node.

        Args:
            child: Child node to remove.

        Raises:
            ValueError: If the node is not a child of this node.
        """
        self.children.remove(child)
        child.parent = None
\end{verbatim}
\end{tcolorbox}
Esta es una clase que hereda los atributos y métodos de la clase ABC que se utiliza para clases abstractas. {\scriptsize \texttt{ABC}} es una clase de la biblioteca estándar utilizada para crear clases abstractas. En esta colección se estudiará más adelante con detalle, por lo que, por ahora, basta con saber que {\scriptsize \texttt{Node}} hereda de ella.
En este ejemplo pueden identificarse fácilmente los principales elementos de una clase. La declaración {\scriptsize \texttt{class Node(ABC):}} indica el nombre de la clase y que hereda de {\scriptsize \texttt{ABC}}. El método {\scriptsize \texttt{\_\_init\_\_}} es el constructor, encargado de inicializar los atributos {\scriptsize \texttt{id}}, {\scriptsize \texttt{parent}} y {\scriptsize \texttt{children}}. Finalmente, los métodos {\scriptsize \texttt{add\_child()}} y {\scriptsize \texttt{remove\_child()}} implementan parte del comportamiento del objeto, permitiendo modificar la estructura del árbol.

\newpage
Para que quede más claro, el esquema sería:


\begin{tcolorbox}
\scriptsize
 \begin{verbatim}
class Node(ABC)
↑
La clase hereda de ABC

↓

__init__
↑
Constructor

↓

self.id
self.parent
self.children
↑
Atributos

↓

add_child
↑
Método
\end{verbatim}
\end{tcolorbox}
\subsection{Conclusión}
Las clases permiten definir nuevos tipos de objetos adaptados a las necesidades del programa.
Estas clases se instancian en objetos que son concreciones de la clase y que representan una instancia concreta del tipo definido por la clase.
Tienen una estructura propia que las hace reconocibles, pero el nombre es a elección del programador y se recomienda usar mayúscula en la primera del nombre explicativo.

\newpage

\section{Atributos}
\subsection{Introducción}
En programación orientada a objetos (POO) las clases se componen principalmente de atributos y métodos. Los atributos son variables asociadas a una clase o a una instancia que almacenan información sobre su estado o sus características.
Los atributos de instancia suelen inicializarse en el método especial {\scriptsize \texttt{\_\_init\_\_}}, aunque también pueden añadirse posteriormente mediante asignación. {\scriptsize \texttt{\_\_init\_\_}} es el método especial encargado de inicializar una instancia recién creada, mientras que la creación del objeto corresponde a {\scriptsize \texttt{\_\_new\_\_}}.
Hay dos tipos principales de atributos: atributos de instancia y atributos de clase. Los de instancia son propios de cada objeto y normalmente obtienen su valor durante la instanciación. Por su parte, los atributos de clase son los que comparten todos los objetos y que pertenecen a la clase.
\subsection{Uso de los atributos}
La definición de los atributos de instancia y de clase se realiza de forma diferente. Para los atributos de instancia se accede al propio objeto, habitualmente mediante {\scriptsize \texttt{self}}, continuado por un punto y el nombre del atributo y asignando un valor. Mientras que para los atributos de clase basta con nombrar el atributo y asignarle un valor.
Una vez que tienen un valor asignado se pueden modificar accediendo a la propiedad y asignando el nuevo valor de la forma: {\scriptsize \texttt{objeto.atributo = valor\_asignado}}.
Tanto las clases como las instancias disponen del atributo especial {\scriptsize \texttt{\_\_dict\_\_}}, que almacena sus atributos en forma de diccionario. Python utiliza estos diccionarios durante la búsqueda de atributos.
Python busca primero un atributo en el {\scriptsize \texttt{\_\_dict\_\_}} de la instancia. Si no lo encuentra, continúa la búsqueda en el {\scriptsize \texttt{\_\_dict\_\_}} de la clase y, posteriormente, en las clases padre si existen.
El shadowing consiste en crear un atributo de instancia con el mismo nombre que un atributo de clase. De este modo, al acceder al atributo desde esa instancia, Python utiliza el atributo de instancia y oculta el de clase.
\subsection{Ejemplo}
El siguiente ejemplo no es parte del código del proyecto. Se ha puesto para entender los conceptos en atributos:


\begin{tcolorbox}
\scriptsize
 \begin{verbatim}
class Circulo:
    
    pi = 3.141592
    
    def __init__(self, radio: float) -> None:
        self.radio: float = radio
        self.diametro: float = 2*radio

\end{verbatim}
\end{tcolorbox}
En el ejemplo, pi es un atributo de clase y no hace falta asignarle un valor cuando se instancia la clase.
Los atributos de radio y diametro se asignan cuando se instancia un objeto.


\begin{tcolorbox}
\scriptsize
 \begin{verbatim}
circulo_1 = Circulo(10.0)

print(circulo_1.diametro)

\end{verbatim}
\end{tcolorbox}
En este fragmento el resultado impreso sería {\scriptsize \texttt{20.0}} pues se pide que imprima el valor del atributo diametro de la instancia {\scriptsize \texttt{circulo\_1}}.
En caso de seguir con el ejemplo anterior:


\begin{tcolorbox}
\scriptsize
 \begin{verbatim}
circulo_1.diametro = 15

print(circulo_1.radio)
\end{verbatim}
\end{tcolorbox}
Esto reasigna el atributo diametro con el valor entero {\scriptsize \texttt{15}}. Sin embargo, no modifica el valor de radio, ya que ambos atributos almacenan valores independientes. Si se desea mantener la relación diametro = 2 * radio, debe implementarse explícitamente, por ejemplo mediante una propiedad ({\scriptsize \texttt{@property}}) o actualizando ambos atributos.
Al asignar {\scriptsize \texttt{circulo\_1.pi = 3}} no se modifica el atributo de clase, sino que se crea un nuevo atributo de instancia llamado {\scriptsize \texttt{pi}}, que oculta al atributo de clase durante la búsqueda de atributos.


\begin{tcolorbox}
\scriptsize
 \begin{verbatim}
circulo_1.pi = 3

circulo_2 = Circulo(5.0)

print(circulo_2.pi)
\end{verbatim}
\end{tcolorbox}
En este ejemplo, se imprimiría 3.141592 pues sólo se ha modificado {\scriptsize \texttt{pi}} para el {\scriptsize \texttt{circulo\_1}}. Esta situación representa lo que es el shadowing.
\subsection{Conclusión}
Los atributos de las clases pueden ser de instancia o de clase.
Los atributos de instancia suelen inicializarse mediante el constructor, aunque también pueden crearse posteriormente mediante asignación.
A los de clase se les asigna desde la clase el valor, que puede llamarse desde cualquier objeto de la clase.
Cuando se accede a un atributo desde una instancia, Python busca primero en los atributos propios del objeto y, si no lo encuentra, continúa la búsqueda en la clase. Este mecanismo explica el comportamiento de los atributos de clase y el fenómeno conocido como shadowing.

\newpage

\section{Métodos}
\subsection{Introducción}
Los métodos son funciones definidas dentro de una clase que describen el comportamiento de sus objetos o de la propia clase. Los métodos de instancia se invocan normalmente desde un objeto previamente instanciado.
La diferencia clave entre funciones y métodos es que los métodos reciben automáticamente una referencia a la instancia ({\scriptsize \texttt{self}}) o a la clase ({\scriptsize \texttt{cls}}), lo que les permite acceder a sus atributos si lo necesitan.
El parámetro {\scriptsize \texttt{self}} representa la instancia sobre la que se está ejecutando el método. Gracias a él es posible acceder a los atributos y métodos del propio objeto.
Hay tres tipos de métodos: de instancia, de clase y estáticos. Los métodos de instancia acceden a los atributos y métodos de la instancia mediante {\scriptsize \texttt{self}}. Por su parte los métodos de clase llaman a los atributos de clase mediante el parámetro {\scriptsize \texttt{cls}}. Finalmente, los métodos estáticos no emplean atributos de la clase, simplemente implementan una operación que utiliza únicamente los parámetros recibidos. En este último caso el método suele estar relacionado con la clase.
Estos métodos no poseen modificadores de acceso {\scriptsize \texttt{public}}, {\scriptsize \texttt{private}} o {\scriptsize \texttt{protected}}. En su lugar utiliza convenciones. Para los métodos públicos se escribe directamente el nombre del método. Los protegidos se representan con {\scriptsize \texttt{\_}} pero no afecta a la visibilidad, sólo al nombre y es una convención que se ha quedado para indicar que es protegido. Los métodos privados usan el doble guion bajo ({\scriptsize \texttt{\_\_}}). Python modifica internamente su nombre (name mangling), lo que evita colisiones con métodos de clases base durante la herencia y dificulta su acceso accidental.
\subsection{Uso de los métodos}
Para invocar un método se escribe: {\scriptsize \texttt{nombre\_objeto.metodo(argumentos)}}. Los métodos pueden modificar el estado del objeto asignando nuevos valores a sus atributos de instancia. También pueden modificar atributos de clase cuando sea necesario.
Son importantes las buenas prácticas a la hora de programar los métodos:

\begin{itemize}
\item  Utilizar nombres que describan acciones ({\scriptsize \texttt{guardar()}}, {\scriptsize \texttt{calcular()}}, {\scriptsize \texttt{abrir()}}).
\item  Cada método debería tener una única responsabilidad.
\item  Evitar métodos excesivamente largos.
\item  Los métodos públicos forman parte de la interfaz de la clase y deberían mantenerse estables.
\item  Reservar los métodos con {\scriptsize \texttt{\_}} para detalles internos de implementación.
\item  Utilizar {\scriptsize \texttt{@staticmethod}} únicamente cuando la función no necesite acceder ni a la instancia ni a la clase.
\item  Utilizar {\scriptsize \texttt{@classmethod}} cuando la operación afecte a la clase o actúe como constructor alternativo.
\end{itemize}
\subsection{Ejemplo}
Para los ejemplos, se va a usar el ejemplo empleado en el apartado de atributos pero añadiendo un método área:

\newpage


\begin{tcolorbox}
\scriptsize
 \begin{verbatim}
class Circulo:
    
    pi = 3.141592
    
    def __init__(self, radio: float) -> None:
        self.radio: float = radio
        self.diametro: float = 2*radio

    def area(self):
        return self.pi * (self.radio**2)

    @classmethod
    def obtener_pi(cls):
        return cls.pi

    @staticmethod
    def suma_areas(area1: float, area2: float) -> float:
        return area1 + area2
\end{verbatim}
\end{tcolorbox}
Como se puede ver en el ejemplo, area depende del atributo de clase pi y del atributo de instancia radio. El método de clase devuelve pi, que es un atributo de clase. Por último, la clase estática suma\_areas devuelve la suma de dos áreas que no se encuentran en ninguno de los atributos de la clase.
\subsection{Conclusión}
Los métodos son funciones propias de las clases que pueden usar los atributos para hacer modificaciones internas u obtener una respuesta en base a los atributos.
Pueden ser de instancia, de clase y estáticos y siguen una convención concreta para métodos públicos, privados y protegidos. En Python no existe un mecanismo estricto de control de acceso como en otros lenguajes. La visibilidad de los métodos se basa principalmente en convenciones de nomenclatura.
Es recomendable seguir buenas prácticas para evitar confusiones y duplicidades.

\newpage


\chapter{Diseño orientado a objetos}

\section{Abstracción}
\subsection{Introducción}
La abstracción es uno de los pilares fundamentales de la Programación Orientada a Objetos (POO). Consiste en representar una entidad del problema mediante una clase que recoge únicamente las características y comportamientos relevantes para el contexto, ignorando los detalles que no aportan información útil.
En otras palabras, una abstracción es un modelo simplificado de un objeto real o conceptual.
La abstracción permite construir modelos más sencillos y fáciles de comprender. Al centrarse únicamente en los aspectos importantes de cada entidad, el código resulta más claro, reutilizable y sencillo de mantener.
Además, facilita el diseño de programas complejos, ya que permite dividir un problema en objetos que representan conceptos bien definidos y con responsabilidades concretas.
\subsection{Uso de la abstracción}
La abstracción comienza identificando las entidades que forman parte del problema que se quiere resolver.
Una vez identificadas, se determina qué información debe almacenar cada una y qué operaciones debe realizar. Esa información se convierte en los atributos de la clase y esas operaciones en sus métodos.
No existe una única abstracción correcta. La forma de modelar una entidad depende del problema que se desea resolver y de la información que resulte relevante en cada caso. Dependiendo del objetivo del programa, una misma entidad puede representarse mediante clases diferentes.
Por ejemplo, una clase Libro puede contener el título, el autor y el ISBN, ya que son características relevantes para una biblioteca. Sin embargo, no es necesario representar el número de fibras del papel o el color exacto de la tinta, porque esa información no resulta útil para ese contexto.
\subsection{Ejemplo}
En este capítulo no es necesario comprender todos los detalles del código; basta con observar cómo la clase representa las características comunes a todos los nodos del árbol.
En Odix todos los elementos del árbol sintáctico (AST) se representan mediante nodos. Aunque existen distintos tipos de nodos, todos comparten unas características comunes: poseen un identificador, conocen su nodo padre y almacenan sus nodos hijos.
Esta idea se representa mediante la siguiente abstracción:


\begin{tcolorbox}
\scriptsize
 \begin{verbatim}
class Node(ABC):
    def __init__(self, node_id: str) -> None:
        self.id = node_id
        self.parent: Node | None = None
        self.children: list[Node] = []
\end{verbatim}
\end{tcolorbox}
La clase {\scriptsize \texttt{Node}} no representa un párrafo, una tabla o un encabezado concreto. Representa el concepto general de nodo de un árbol sintáctico, definiendo únicamente las características comunes a todos ellos.
Posteriormente, otras clases derivadas heredarán de ella para añadir el comportamiento específico de cada tipo de nodo.
\subsection{Conclusión}
La abstracción consiste en construir modelos simplificados de las entidades del problema, conservando únicamente la información y el comportamiento necesarios para el contexto. Suele ser el primer paso en el diseño orientado a objetos, ya que permite identificar las clases que formarán parte del programa y definir sus responsabilidades antes de comenzar su implementación.

\newpage

\section{Herencia}
\subsection{Introducción}
Uno de los conceptos más importantes de la programación orientada a objetos (POO) es la herencia. La herencia es un mecanismo de la programación orientada a objetos que permite que una clase derive de otra, heredando sus atributos y métodos y pudiendo ampliarlos o modificarlos.
La clase padre es la clase inicial que define los atributos y métodos comunes que tienen todas las clases hijas. La clase hija puede reutilizar los atributos y métodos heredados, añadir nuevos y sobrescribir aquellos métodos cuyo comportamiento necesite especializar.
Una clase hija representa un caso más específico de la clase padre. Mantiene todas las características generales de ésta y añade únicamente el comportamiento o los datos que la diferencian.
La función {\scriptsize \texttt{super()}} permite acceder a los métodos de la clase padre desde la clase hija. Se utiliza habitualmente para reutilizar parte de su implementación, por ejemplo llamando al constructor de la clase padre antes de inicializar los nuevos atributos de la clase hija.
Las principales ventajas de la herencia son la reutilización de código, la reducción de duplicidades, la mejora del mantenimiento y la posibilidad de crear especializaciones a partir de una misma clase base.
\subsection{Uso de la herencia}
La herencia en Python es sencilla y se realiza creando la clase padre y una vez creada, añadiéndola entre paréntesis justo después del nombre de la clase: {\scriptsize \texttt{class Nombre\_hija(Nombre\_padre):}}.
Toda clase hija también puede utilizarse allí donde se espere un objeto de la clase padre. Con un ejemplo, una llave allen es un tipo de herramienta, un destornillador es otra herramienta. Se le puede decir a una función que la entrada son llaves allen o destornilladores, pero eso obligaría a tener una función por cada una. Si la función requiere de una herramienta, se le puede pasar un objeto tipo herramienta a la función, permitiendo utilizar ambas clases para la función con un único bloque de código.
\subsection{Ejemplo}
La clase {\scriptsize \texttt{Herramienta}} contendrá el comportamiento común, mientras que {\scriptsize \texttt{Allen}} y {\scriptsize \texttt{Destornillador}} representarán dos especializaciones diferentes.

\newpage


\begin{tcolorbox}
\scriptsize
 \begin{verbatim}
class Herramienta:

    tipo = "herramienta"
    
    def __init__(self, size:float) -> None:
        self.size = size

    def ensamblar(self) -> None:
        pass
\end{verbatim}
\end{tcolorbox}
Definimos dos clases hijas nuevas de la clase Herramienta: Las clases Allen y Destornillador:


\begin{tcolorbox}
\scriptsize
 \begin{verbatim}
class Allen(Herramienta):

    def __init__(self, size:float, length:float) -> None:
        super().__init__(size)
        self.length = length

    def ensamblar(self) -> None:
        print("Apretar con llave Allen")
\end{verbatim}
\end{tcolorbox}


\begin{tcolorbox}
\scriptsize
 \begin{verbatim}
class Destornillador(Herramienta):

    def __init__(self, size:float, punta:str) -> None:
        super().__init__(size)
        self.punta = punta

    def ensamblar(self):
        print("Apretar con llave Destornillador")
\end{verbatim}
\end{tcolorbox}
Ahora, al iniciar un Destornillador necesitamos:


\begin{tcolorbox}
\scriptsize
 \begin{verbatim}
destornillador_1 = Destornillador(size = 8.0, punta = "Estrella")
\end{verbatim}
\end{tcolorbox}
Gracias a la herencia, tanto la clase {\scriptsize \texttt{Allen}} como {\scriptsize \texttt{Destornillador}} reutilizan el atributo {\scriptsize \texttt{size}} definido en {\scriptsize \texttt{Herramienta}}, añadiendo únicamente los atributos que las diferencian ({\scriptsize \texttt{length}} y {\scriptsize \texttt{punta}}, respectivamente). De este modo se evita duplicar código y se obtiene una estructura más modular y fácil de mantener.
\subsection{Conclusión}
La herencia permite reutilizar código común entre varias clases relacionadas, facilitando el mantenimiento y permitiendo crear especializaciones de una misma clase base sin duplicar comportamiento.

\newpage

\section{Composición}
\subsection{Introducción}
La composición es un mecanismo mediante el cual una clase utiliza objetos de otras clases para construir su comportamiento. 
Contrasta con la herencia en que al heredar las clases hijas ganan todos los atributos y métodos de la clase padre aunque luego puedan sobreescribirlos. En la composición la clase no hereda los atributos ni los métodos de las clases que la forman, sino que puede utilizarlos a través de los objetos que contiene. La herencia modela una relación de especialización ("es un"), mientras que la composición modela una relación de agregación ("tiene un").
Se suele preferir la composición frente a la herencia a la hora de programar porque tiene un menor acoplamiento y ofrece una mayor flexibilidad. De ahí la conocida recomendación: \textit{"Favor composition over inheritance."}
\subsection{Uso de la composición}
Para usar la composición basta con almacenar objetos de otras clases como atributos de una clase. Estos objetos pueden crearse internamente o recibirse como parámetros durante la inicialización.
Una vez que se tiene la instancia dentro de la clase, esta puede utilizar los métodos y atributos de los objetos que la componen.
\subsection{Ejemplo}
En el siguiente ejemplo se ve cómo es una composición:


\begin{tcolorbox}
\scriptsize
 \begin{verbatim}
class EmailService:

    def enviar(self, mensaje):
        print(f"Email: {mensaje}")


class Persona:

    def __init__(self):
        self.email = EmailService()

    def avisar(self):
        self.email.enviar("Hola")
\end{verbatim}
\end{tcolorbox}
 
La clase {\scriptsize \texttt{Persona}} contiene un objeto de la clase {\scriptsize \texttt{EmailService}}, cuyo método {\scriptsize \texttt{enviar()}} utiliza para enviar un mensaje. 
Por lo tanto, {\scriptsize \texttt{Persona}} no hereda los métodos de {\scriptsize \texttt{EmailService}}, sino que los utiliza a través del objeto {\scriptsize \texttt{email}}. {\scriptsize \texttt{Persona}} delega la responsabilidad de enviar correos en {\scriptsize \texttt{EmailService}}.
\subsection{Conclusión}
En POO la composición permite construir objetos complejos a partir de otros más simples, delegando responsabilidades entre ellos y favoreciendo un diseño más flexible y con menor acoplamiento que la herencia. Por ello, suele ser la opción preferida frente a la herencia cuando ambas permiten resolver el mismo problema.

\newpage

\section{Encapsulación}
\subsection{Introducción}
La encapsulación consiste en diseñar una clase de forma que el acceso y la modificación de su estado se realicen a través de una interfaz pública controlada, ocultando por convención los detalles de implementación interna.
En Python, la encapsulación suele implementarse mediante propiedades y otros métodos públicos que controlan el acceso y la modificación del estado interno del objeto, mientras que los atributos internos se consideran detalles de implementación.
\subsection{Uso de la encapsulación}
Para implementar la encapsulación, Python se apoya principalmente en convenciones de nomenclatura y en propiedades ({\scriptsize \texttt{property}}) que permiten controlar el acceso al estado interno del objeto. En muchos lenguajes de programación, la encapsulación impide acceder directamente a los atributos privados o protegidos, obligando a utilizar métodos públicos, habitualmente denominados \textit{setters} y \textit{getters}, para modificar o consultar su valor.
La convención para atributos protegidos es el uso de guión bajo ({\scriptsize \texttt{\_}}) y para privados doble guión bajo ({\scriptsize \texttt{\_\_}}) antecediendo el nombre del atributo. Estas convenciones no impiden el poder modificar el valor asignado al atributo pero en Python se sigue la filosofía: \textit{"We're all consenting adults here"}.
La palabra clave {\scriptsize \texttt{@property}} permite exponer un método como si fuera un atributo. Si además se define un método decorado con {\scriptsize \texttt{@atributo.setter}}, la asignación sobre dicho atributo ejecutará automáticamente ese método. 
La encapsulación ofrece varias ventajas:

\begin{itemize}
\item  Permite validar los datos antes de modificar el estado interno.
\item  Reduce el acoplamiento ocultando los detalles de implementación.
\item  Facilita modificar la implementación interna sin afectar al código que utiliza la clase.
\end{itemize}
\subsection{Ejemplo}
Para el ejemplo de encapsulación se empleará un ejemplo con la clase {\scriptsize \texttt{Circulo}}:

\newpage


\begin{tcolorbox}
\scriptsize
 \begin{verbatim}
class Circulo:

    def __init__(self, radio:float):
        self.__radio:float = radio

    @property
    def radio(self) -> float:
        return self.__radio

    @radio.setter
    def radio(self, nuevo_radio:float) -> None:
        if nuevo_radio > 0.0:
            self.__radio = nuevo_radio

    @property
    def diametro(self) -> float:
        return self.__radio * 2

    
\end{verbatim}
\end{tcolorbox}
En este ejemplo {\scriptsize \texttt{\_\_radio}} es un atributo privado. La propiedad radio permite consultar su valor y el método decorado con {\scriptsize \texttt{@radio.setter}} controla su modificación, validando que el nuevo radio sea mayor que cero.
Para invocar cada uno se puede ejecutar:


\begin{tcolorbox}
\scriptsize
 \begin{verbatim}
circulo_1 = Circulo(5.0)

print(circulo_1.radio)

circulo_1.radio = 4.0

print(circulo_1.diametro)
\end{verbatim}
\end{tcolorbox}
En el primer print se consulta la propiedad lo que lleva a obtener el valor del radio ({\scriptsize \texttt{5.0}} en este punto). En el segundo print se obtiene el valor del diámetro que es calculado en base al radio. Sin embargo, este diámetro ya no es el doble de {\scriptsize \texttt{5.0}}, sino que es {\scriptsize \texttt{8.0}} ya que se invocó el setter justo antes del print para modificar el radio a {\scriptsize \texttt{4.0}}.
\subsection{Conclusión}
La encapsulación en Python es más un convenio que una característica intrínseca del lenguaje. No existe la ocultación de atributos como tal pero se pueden usar convenciones para poder usar propiedades y métodos \textit{setter} como en otros lenguajes de programación.
Esto permite validar, transformar o calcular valores antes de modificar el estado interno del objeto, además de mantener una interfaz pública estable aunque cambie la implementación interna de la clase.

\newpage

\section{Polimorfismo}
\subsection{Introducción}
En programación orientada a objetos (POO) se define polimorfismo como la capacidad de utilizar una misma interfaz para trabajar con objetos de diferentes tipos, es decir, una misma función puede trabajar con objetos de diferentes clases utilizando la misma interfaz, sin necesidad de conocer el tipo concreto del objeto.
En Python el polimorfismo puede conseguirse principalmente de dos formas: mediante herencia o mediante Duck Typing. El polimorfismo mediante herencia emplea la sobreescritura de un método de la clase padre en las clases hijas, haciendo que todos los objetos de ese tipo puedan ser empleados en la misma función. El polimorfismo sin herencia se basa en el \textbf{Duck Typing}: "Si camina como un pato y hace cuac como un pato, entonces lo tratamos como un pato.". En este caso no es necesario que exista una relación de herencia entre las clases. Basta con que el objeto implemente el método esperado para que pueda utilizarse.
El polimorfismo permite evitar duplicidades de código y simplificar el uso de los objetos, ya que una misma función puede invocar un método del objeto sin necesidad de conocer la clase a la que pertenece. 
\subsection{Uso del polimorfismo}
Para emplear el polimorfismo se puede hacer uso de la herencia o no. Para llevar a cabo un polimorfismo mediante herencia, la clase padre tendrá que definir uno o varios métodos que en las clases hijas, especializaciones de la clase padre, sobreescribirán pero manteniendo el mismo nombre del método.
Cuando se quiere realizar polimorfismo sin herencia, basta con incluir el método con las mismas propiedades dentro de una clase ya que después la función no diferencia que un objeto sea de una clase o no, simplemente intenta acceder al método solicitado. Si el objeto lo implementa, la llamada se realiza correctamente; en caso contrario se producirá una excepción ({\scriptsize \texttt{AttributeError}}).
Por lo tanto, para implementar el polimorfismo en Python, es suficiente con el uso de métodos con el mismo nombre y argumentos dentro de distintas clases que luego, al instanciarlas en objetos, podrán ser utilizadas por las funciones independientemente de si son de la misma clase o clase padre que tengan.
\subsection{Ejemplo}
Para este ejemplo que se presentó en los métodos de la clase, la clase {\scriptsize \texttt{Herramienta}} contendrá el comportamiento común, mientras que {\scriptsize \texttt{Allen}} y {\scriptsize \texttt{Destornillador}} representarán dos especializaciones diferentes.

\newpage


\begin{tcolorbox}
\scriptsize
 \begin{verbatim}
class Herramienta:

    tipo = "herramienta"
    
    def __init__(self, size:float) -> None:
        self.size = size

    def ensamblar(self) -> None:
        raise NotImplementedError
\end{verbatim}
\end{tcolorbox}
Definimos dos clases hijas nuevas de la clase Herramienta: Las clases Allen y Destornillador:


\begin{tcolorbox}
\scriptsize
 \begin{verbatim}
class Allen(Herramienta):

    def __init__(self, size:float, length:float) -> None:
        super().__init__(size)
        self.length = length

    def ensamblar(self) -> None:
        print("Apretar con llave Allen")
\end{verbatim}
\end{tcolorbox}


\begin{tcolorbox}
\scriptsize
 \begin{verbatim}
class Destornillador(Herramienta):

    def __init__(self, size:float, punta:str) -> None:
        super().__init__(size)
        self.punta = punta

    def ensamblar(self):
        print("Apretar con llave Destornillador")
\end{verbatim}
\end{tcolorbox}
Como se puede ver, {\scriptsize \texttt{Allen}} y {\scriptsize \texttt{Destornillador}} sobrescriben el método ensamblar heredado de {\scriptsize \texttt{Herramienta}}. Si definimos una función ensamble, podemos usar las herramientas llamando al método ensamblar:


\begin{tcolorbox}
\scriptsize
 \begin{verbatim}
def ensamble(tool:Herramienta) -> None:
    tool.ensamblar()
\end{verbatim}
\end{tcolorbox}
Según la herramienta que se ponga usará el método de esa herramienta. Python no necesita comprobar si el objeto es una instancia de {\scriptsize \texttt{Allen}} o {\scriptsize \texttt{Destornillador}}; simplemente intenta acceder al método {\scriptsize \texttt{ensamblar()}}.
Aunque el parámetro se haya anotado como {\scriptsize \texttt{Herramienta}}, esta anotación no impide pasar objetos de otras clases. Los \textit{type hints} sirven como ayuda para el programador y para las herramientas de análisis estático, pero no modifican el comportamiento del programa en tiempo de ejecución.
Un ejemplo de Duck Typing sería el siguiente:


\begin{tcolorbox}
\scriptsize
 \begin{verbatim}

class Robot:

    def ensamblar(self):
        print("Ensamblando automáticamente")

robot = Robot()

ensamble(robot)
\end{verbatim}
\end{tcolorbox}
Aunque Robot no hereda de {\scriptsize \texttt{Herramienta}}, la función sigue funcionando porque el objeto implementa el método {\scriptsize \texttt{ensamblar()}}.
\subsection{Conclusión}
El polimorfismo permite emplear objetos de distintas clases en la misma función siempre que implementen la interfaz esperada (normalmente uno o varios métodos con el mismo nombre y comportamiento compatible). Esto evita tener que hacer casos concretos para cada clase que se le pasa al método.
Es un concepto de POO muy empleado que en Python se puede hacer con herencia o sin herencia.
El polimorfismo es uno de los pilares de la programación orientada a objetos, ya que permite escribir funciones más genéricas, reutilizables y fáciles de mantener.

\newpage


\chapter{Herramientas de Python}

\section{Clases abstractas}
\subsection{Introducción}
Una clase abstracta es una clase que sirve como base para otras clases y que no está pensada para ser instanciada directamente. En Python las clases abstractas se implementan mediante las Abstract Base Class (ABC) proporcionadas por el módulo {\scriptsize \texttt{abc}}.
Una clase representa un concepto mediante sus propiedades y sus métodos. La abstracción consiste en quedarse con las características esenciales de un objeto y ocultar los detalles que no importan. En el proceso de abstracción se identifican las propiedades y métodos esenciales.
Cuando varias clases comparten comportamiento o representan variaciones de un mismo concepto, es habitual crear una clase padre de la que todas hereden las propiedades y métodos de esta.
Este enfoque permite modularizar y reutilizar el código dando como resultado un código más mantenible y flexible. Por lo tanto la herencia es una herramienta muy útil cuando existe una relación clara entre las clases.
Sin embargo, la herencia por sí sola no hace que las clases hijas implementen determinados métodos o propiedades. Es posible no definir alguno y que el error aparezca mucho más adelante durante la ejecución. Para evitar este problema existen las Abstract Base Class (ABC).
\subsection{Uso de las ABC}
Una clase abstracta es una clase que sirve como base para otras clases y que puede obligar a sus clases derivadas a implementar determinados métodos o propiedades. En Python este comportamiento se implementa heredando de {\scriptsize \texttt{ABC}}. En el caso de no definir lo que la clase base indica que debe ser definido, se manda un mensaje de error y se para la ejecución del programa.
Dicho en otras palabras, la clase base obliga a sus hijas a seguir una guía que marca esa clase base. Es una garantía de que todas las clases derivadas implementan la interfaz definida por la clase abstracta. Así se puede asegurar que cada vez que se defina una clase hija de una de éstas clases todas sigan un patrón común.
Una vez definidas las clases hijas, se pueden crear funciones y métodos que reciban cualquier objeto de la clase base. Eso permite incluir todos los hijos de la clase base como entrada en esas funciones.
En caso de olvidarse de incluir alguna de las propiedades o métodos lanzará una excepción (TypeError) que impide crear el objeto, evitando que el programa continúe con una implementación incompleta. De este modo se garantiza que todas las clases derivadas implementan la interfaz definida por la clase abstracta, evitando implementaciones incompletas y errores difíciles de detectar durante la ejecución.
\subsection{Ejemplo}
Para crear una clase abstracta basta con heredar de ABC, una clase incluida en el módulo abc de Python. 
Vamos con un ejemplo mecánico: Definimos una clase {\scriptsize \texttt{Pieza}} que representa en abstracto a cualquier pieza. La clase {\scriptsize \texttt{Pieza}} define la interfaz común que deberán implementar todas las piezas concretas.
Así se marcan las propiedades que deben tener las piezas con su tamaño y la herramienta que se usa y también el método de ensamblaje:

\newpage


\begin{tcolorbox}
\scriptsize
 \begin{verbatim}

from abc import ABC, abstractmethod

class Pieza(ABC):

    @property
    @abstractmethod
    def tamaño(self):
        pass

    @property
    @abstractmethod
    def herramienta(self):
        pass

    @abstractmethod
    def ensamblar(self):
        pass
\end{verbatim}
\end{tcolorbox}
Definimos dos clases hijas nuevas de la clase {\scriptsize \texttt{Pieza}}: Las clases TuercaM8 y TornilloM8:

\newpage


\begin{tcolorbox}
\scriptsize
 \begin{verbatim}
class TuercaM8(Pieza):

    @property
    def tamaño(self):
        return "M8"

    @property
    def herramienta(self):
        return "Llave inglesa de 13 mm"

    def ensamblar(self):
        print("Apretar a 50 N·m")
\end{verbatim}
\end{tcolorbox}


\begin{tcolorbox}
\scriptsize
 \begin{verbatim}
class TornilloM8(Pieza):

    @property
    def tamaño(self):
        return "M8"

    @property
    def herramienta(self):
        return "Llave Allen"

    def ensamblar(self):
        print("Apretar con llave Allen")
\end{verbatim}
\end{tcolorbox}
Estas clases están completas porque definen todo lo que la clase base marca. Una clase incompleta sería:

\newpage


\begin{tcolorbox}
\scriptsize
 \begin{verbatim}
class Remache(Pieza):

    @property
    def tamaño(self):
        return "4 mm"
\end{verbatim}
\end{tcolorbox}
A la que le faltan la propiedad herramienta y el método ensamblar. Para comprobar el funcionamiento se puede ejecutar:


\begin{tcolorbox}
\scriptsize
 \begin{verbatim}
remache = Remache()
\end{verbatim}
\end{tcolorbox}
Lo cual lanzaría un error:


\begin{tcolorbox}
\scriptsize
 \begin{verbatim}
TypeError:
Can't instantiate abstract class Remache
with abstract methods herramienta, ensamblar
\end{verbatim}
\end{tcolorbox}
Como se puede ver, en caso de no cumplir las condiciones de la clase padre, no se puede seguir con el desarrollo.
Por otro lado, para usar en una función las clases no es necesario usar cada clase por separado como en este ejemplo:


\begin{tcolorbox}
\scriptsize
 \begin{verbatim}
def montar(tuerca: TuercaM8):
    ...
\end{verbatim}
\end{tcolorbox}
Se puede hacer directamente una llamada a la clase base:


\begin{tcolorbox}
\scriptsize
 \begin{verbatim}
def montar(pieza: Pieza):
    pieza.ensamblar()
\end{verbatim}
\end{tcolorbox}
Una vez que todas las clases cumplen la interfaz definida por la clase abstracta, pueden utilizarse indistintamente allí donde se espere un objeto del tipo Pieza. Este comportamiento se estudiará con detalle en la píldora dedicada al polimorfismo.
\subsection{Conclusión}
Las clases abstractas no están pensadas para ser instanciadas directamente, sino para servir de base a otras clases.
Las ABC permiten definir unas pautas en el diseño. Convierten las decisiones de diseño en reglas que el propio lenguaje hace cumplir, reduciendo errores y facilitando el mantenimiento del código.

\newpage

\section{Métodos especiales}
\subsection{Introducción}
Existen determinadas operaciones del lenguaje Python cuyo comportamiento está definido por el propio intérprete. Por ejemplo, cuando se evalúa la expresión {\scriptsize \texttt{a + b}}, Python debe decidir qué operación realizar sobre ambos objetos. Los métodos que permiten definir cómo debe comportarse un objeto cuando participa en estas operaciones reciben el nombre de métodos especiales o \textit{dunder methods} (de double underscore methods). También se conocen como métodos mágicos (magic methods), aunque el término dunder methods es el más utilizado actualmente.
Los métodos especiales permiten que los objetos personalizados se integren con el propio lenguaje Python. Gracias a ellos, un objeto puede comportarse como una cadena, una lista, un número o cualquier otro tipo compatible con las operaciones del lenguaje.
Gracias a los métodos especiales es posible:

\begin{itemize}
\item  Personalizar la representación textual de un objeto.
\item  Comparar objetos entre sí.
\item  Sobrecargar operadores como {\scriptsize \texttt{+}}, {\scriptsize \texttt{-}} o {\scriptsize \texttt{*}}.
\item  Hacer que un objeto sea iterable.
\item  Permitir que pueda utilizarse con {\scriptsize \texttt{len()}}.
\item  Convertir un objeto en una función invocable.
\item  Gestionar el acceso a atributos.
\item  Implementar gestores de contexto ({\scriptsize \texttt{with}}).
\end{itemize}
En lugar de crear funciones específicas para cada una de estas operaciones, Python utiliza un conjunto estandarizado de métodos especiales que el intérprete invoca automáticamente cuando son necesarios.
\subsection{Uso de los dunder methods}
Los métodos especiales no suelen llamarse directamente. Cuando se emplean determinadas funciones integradas de Python, como {\scriptsize \texttt{len()}} o {\scriptsize \texttt{str()}}, el intérprete invoca automáticamente el método especial correspondiente.
Algunos métodos especiales son invocados mediante funciones integradas de Python, mientras que otros se ejecutan automáticamente al utilizar operadores o determinadas construcciones del lenguaje:
\textbf{No todos los métodos especiales son llamados por una función de Python}. Algunos son invocados por funciones como {\scriptsize \texttt{len()}} o {\scriptsize \texttt{str()}}, mientras que otros son activados automáticamente por el intérprete al utilizar operadores ({\scriptsize \texttt{+}}, {\scriptsize \texttt{==}}, {\scriptsize \texttt{[]}}) o determinadas construcciones del lenguaje ({\scriptsize \texttt{with}}, acceso a atributos, iteración, etc.).
Buenas prácticas a la hora de usar estos métodos son:

\begin{itemize}
\item  Implementar únicamente los métodos especiales que aporten un comportamiento útil a la clase.
\item  Mantener el significado esperado de cada operación para evitar comportamientos sorprendentes.
\item  Preferir funciones del lenguaje (len(), print(), operadores...) en lugar de llamar directamente a los métodos especiales.
\item  Consultar la documentación oficial cuando se implementen métodos menos habituales, ya que algunos tienen requisitos específicos.
\end{itemize}
\subsection{Ejemplo}
Se toma el ejemplo de la clase Allen, que es una llave, y se usará un método especial para aplicar la operación suma con una clase tornillo y esta herramienta. Se definen las clases:

\newpage


\begin{tcolorbox}
\scriptsize
 \begin{verbatim}
class Herramienta:

    tipo = "herramienta"
    
    def __init__(self, size:float) -> None:
        self.size = size

    def ensamblar(self) -> str:
        raise NotImplementedError

class Pieza:
    tipo = "pieza"

    def __init__(self, size:float) -> None:
        self.size = size

    def ensamblar(self) -> str:
        raise NotImplementedError


class Allen(Herramienta):

    def __init__(self, size:float, length:float) -> None:
        super().__init__(size)
        self.length = length

    def ensamblar(self) -> str:
        return "Apretado con llave Allen"

    def __add__(self,pieza:Pieza) -> str:
        return f"{self.ensamblar()} {pieza.ensamblar()}"

class Tornillo(Pieza):

    def __init__(self, size:float, length:float) -> None:
        super().__init__(size)
        self.length = length

    def ensamblar(self) -> str:
        return "Tornillo"

    
    def __add__(self,hta:Herramienta) -> str:
        return f"{self.ensamblar()} {hta.ensamblar()}"
\end{verbatim}
\end{tcolorbox}
Que después se utilizaría de la siguiente forma:


\begin{tcolorbox}
\scriptsize
 \begin{verbatim}
allen_1 = Allen(8.0, 10.0)

tornillo_1 = Tornillo(8.0, 5.0)

print(allen_1 + tornillo_1)

print(tornillo_1 + allen_1)
\end{verbatim}
\end{tcolorbox}
Al evaluar la expresión {\scriptsize \texttt{allen\_1 + tornillo\_1}}, Python no realiza la suma directamente, sino que invoca automáticamente el método {\scriptsize \texttt{allen\_1.\_\_add\_\_(tornillo\_1)}}.
En este ejemplo el resultado depende del orden de los operandos, ya que cada clase implementa su propia versión de {\scriptsize \texttt{\_\_add\_\_()}}. Esto no supone un error, pero sí muestra que la sobrecarga de operadores debe utilizarse con criterio para que el comportamiento de los objetos resulte intuitivo.
\subsection{Conclusión}
Los métodos especiales permiten integrar las clases personalizadas con el propio funcionamiento del lenguaje Python. Gracias a ellos es posible redefinir cómo se representan los objetos, cómo se comparan, cómo responden a operadores, cómo se recorren o cómo interactúan con funciones integradas como {\scriptsize \texttt{len()}} o {\scriptsize \texttt{str()}}.
Aunque son una herramienta muy potente, conviene implementar únicamente aquellos métodos especiales que aporten un comportamiento claro y coherente, evitando redefinir operaciones de forma que resulten confusas para otros desarrolladores.

\newpage


\chapter{Anexos}

\subsection{Anexo A. Tipos de datos de Python}

\begin{tcolorbox}
\begin{tabular}{ll}
 \textbf{Tipo de dato}   &  \textbf{Clase}    \\
 Número entero  &  {\scriptsize \texttt{int}}    \\
 Decimal        &  {\scriptsize \texttt{float}}  \\
 Cadena         &  {\scriptsize \texttt{str}}    \\
 Booleano       &  {\scriptsize \texttt{bool}}   \\
 Lista          &  {\scriptsize \texttt{list}}   \\
 Tupla          &  {\scriptsize \texttt{tuple}}  \\
 Conjunto       &  {\scriptsize \texttt{set}}    \\
 Diccionario    &  {\scriptsize \texttt{dict}}   \\
\end{tabular}
\end{tcolorbox}

\newpage
\subsection{Anexo B. Métodos más utilizados}
\subsubsection{Listas}

\begin{tcolorbox}
\begin{tabular}{ll}
 \textbf{Método}       &  \textbf{Descripción}                           \\
 {\scriptsize \texttt{append()}}   &  Añade un elemento al final.           \\
 {\scriptsize \texttt{insert()}}   &  Inserta un elemento en una posición.  \\
 {\scriptsize \texttt{remove()}}   &  Elimina un elemento.                  \\
 {\scriptsize \texttt{pop()}}      &  Elimina y devuelve un elemento.       \\
 {\scriptsize \texttt{sort()}}     &  Ordena la lista.                      \\
 {\scriptsize \texttt{reverse()}}  &  Invierte el orden.                    \\
 {\scriptsize \texttt{clear()}}    &  Vacía la lista.                       \\
 {\scriptsize \texttt{index()}}    &  Devuelve el índice en el que se encuentra el elemento  \\
\end{tabular}
\end{tcolorbox}
\subsubsection{Cadenas}

\begin{tcolorbox}
\begin{tabular}{ll}
 \textbf{Método}       &  \textbf{Descripción}              \\
 {\scriptsize \texttt{lower()}}    &  Convierte a minúsculas.  \\
 {\scriptsize \texttt{upper()}}    &  Convierte a mayúsculas.  \\
 {\scriptsize \texttt{strip()}}    &  Elimina espacios.        \\
 {\scriptsize \texttt{replace()}}  &  Sustituye texto.         \\
 {\scriptsize \texttt{split()}}    &  Divide una cadena.       \\
 {\scriptsize \texttt{join()}}     &  Une los elementos de una colección. \\
 {\scriptsize \texttt{startswith()}} &  Comprueba si la cadena comienza por un texto. \\
 {\scriptsize \texttt{endswith()}} &   Comprueba si la cadena termina en un texto. \\
\end{tabular}
\end{tcolorbox}
\subsubsection{Diccionarios}

\begin{tcolorbox}
\begin{tabular}{ll}
 \textbf{Método}      &  \textbf{Descripción}                        \\
 {\scriptsize \texttt{keys()}}    &  Devuelve las claves.               \\
 {\scriptsize \texttt{values()}}  &  Devuelve los valores.              \\
 {\scriptsize \texttt{items()}}   &  Devuelve pares clave-valor.        \\
 {\scriptsize \texttt{get()}}     &  Obtiene un valor de forma segura.  \\
 {\scriptsize \texttt{update()}}  &  Actualiza el diccionario.          \\
 {\scriptsize \texttt{pop()}}     &  Elimina una clave.                 \\
 {\scriptsize \texttt{popitem()}} &  Elimina y devuelve el último par clave-valor. \\
\end{tabular}
\end{tcolorbox}
\subsubsection{Conjuntos}

\begin{tcolorbox}
\begin{tabular}{ll}
 \textbf{Método}            &  \textbf{Descripción}           \\
 {\scriptsize \texttt{add()}}           &  Añade un elemento.    \\
 {\scriptsize \texttt{remove()}}        &  Elimina un elemento.  \\
 {\scriptsize \texttt{union()}}         &  Unión de conjuntos.   \\
 {\scriptsize \texttt{intersection()}}  &  Intersección.         \\
 {\scriptsize \texttt{difference()}}    &  Devuelve la diferencia de los elementos. \\
\end{tabular}
\end{tcolorbox}
\subsection{Anexo C. Operadores}

\begin{tcolorbox}
\begin{tabular}{ll}
 \textbf{Tipo}         &  \textbf{Operadores}         \\
 Aritméticos  &  {\scriptsize \texttt{+ - * / // \% **}}  \\
 Comparación  &  {\scriptsize \texttt{== != \textgreater{} \textless{} \textgreater{}= \textless{}=}}  \\
 Lógicos      &  {\scriptsize \texttt{and or not}}       \\
\end{tabular}
\end{tcolorbox}

\newpage
\subsection{Anexo D. Conversión de tipos}

\begin{tcolorbox}
\begin{tabular}{ll}
 \textbf{Función}    &  \textbf{Convierte a}   \\
 {\scriptsize \texttt{int()}}    &  Entero        \\
 {\scriptsize \texttt{float()}}  &  Decimal       \\
 {\scriptsize \texttt{str()}}    &  Cadena        \\
 {\scriptsize \texttt{bool()}}   &  Booleano      \\
 {\scriptsize \texttt{list()}}   &  Lista         \\
 {\scriptsize \texttt{tuple()}}  &  Tupla         \\
 {\scriptsize \texttt{set()}}    &  Conjunto      \\
 {\scriptsize \texttt{dict()}}   &  Diccionario*  \\
\end{tabular}
\end{tcolorbox}
*No todas las conversiones son posibles. La conversión solo puede realizarse cuando el valor original puede representarse mediante el nuevo tipo de dato.

\newpage
\subsection{Anexo E. Precedencia de operadores}
Cuando una expresión contiene varios operadores, Python los evalúa siguiendo el siguiente orden de prioridad.

\begin{tcolorbox}
\begin{tabular}{ll}
\textbf{Precedencia} & \textbf{Operador} \\
1 & {\scriptsize \texttt{()}} \\
2 & {\scriptsize \texttt{**}} \\
3 & {\scriptsize \texttt{* / // \%}} \\
4 & {\scriptsize \texttt{+ -}} \\
5 & {\scriptsize \texttt{\textless{} \textgreater{} ==}} \\
6 & {\scriptsize \texttt{not}} \\
7 & {\scriptsize \texttt{and}} \\
8 & {\scriptsize \texttt{or}} \\
\end{tabular}
\end{tcolorbox}

\newpage
\subsection{Anexo F. Funciones integradas utilizadas en este volumen}

\begin{tcolorbox}
\begin{tabular}{ll}
 \textbf{Función}    &  \textbf{Descripción}                        \\
 {\scriptsize \texttt{print()}}  &  Muestra información por pantalla.  \\
 {\scriptsize \texttt{type()}}   &  Devuelve el tipo de un valor.      \\
 {\scriptsize \texttt{len()}}    &  Devuelve el número de elementos.   \\
 {\scriptsize \texttt{int()}}    &  Convierte a entero.                \\
 {\scriptsize \texttt{float()}}  &  Convierte a decimal.               \\
 {\scriptsize \texttt{str()}}    &  Convierte a cadena.                \\
 {\scriptsize \texttt{bool()}}   &  Convierte a booleano.              \\
\end{tabular}
\end{tcolorbox}

\newpage
\subsection{Anexo G: Relación entre operaciones de Python y métodos especiales}

\begin{tcolorbox}
\begin{tabular}{ll}
 \textbf{Operación en Python}    &  \textbf{Método especial invocado}              \\
 {\scriptsize \texttt{str(obj)}}             &  {\scriptsize \texttt{obj.\_\_str\_\_()}}                       \\
 {\scriptsize \texttt{repr(obj)}}            &  {\scriptsize \texttt{obj.\_\_repr\_\_()}}                      \\
 {\scriptsize \texttt{format(obj)}}          &  {\scriptsize \texttt{obj.\_\_format\_\_()}}                    \\
 {\scriptsize \texttt{bytes(obj)}}           &  {\scriptsize \texttt{obj.\_\_bytes\_\_()}}                     \\
 {\scriptsize \texttt{len(obj)}}             &  {\scriptsize \texttt{obj.\_\_len\_\_()}}                       \\
 {\scriptsize \texttt{bool(obj)}}            &  {\scriptsize \texttt{obj.\_\_bool\_\_()}}                      \\
 {\scriptsize \texttt{int(obj)}}             &  {\scriptsize \texttt{obj.\_\_int\_\_()}}                       \\
 {\scriptsize \texttt{float(obj)}}           &  {\scriptsize \texttt{obj.\_\_float\_\_()}}                     \\
 {\scriptsize \texttt{complex(obj)}}         &  {\scriptsize \texttt{obj.\_\_complex\_\_()}}                   \\
 {\scriptsize \texttt{hash(obj)}}            &  {\scriptsize \texttt{obj.\_\_hash\_\_()}}                      \\
 {\scriptsize \texttt{iter(obj)}}            &  {\scriptsize \texttt{obj.\_\_iter\_\_()}}                      \\
 {\scriptsize \texttt{next(iterador)}}       &  {\scriptsize \texttt{iterador.\_\_next\_\_()}}                 \\
 {\scriptsize \texttt{reversed(obj)}}        &  {\scriptsize \texttt{obj.\_\_reversed\_\_()}}                  \\
 {\scriptsize \texttt{obj[indice]}}          &  {\scriptsize \texttt{obj.\_\_getitem\_\_(indice)}}             \\
 {\scriptsize \texttt{obj[indice] = valor}}  &  {\scriptsize \texttt{obj.\_\_setitem\_\_(indice, valor)}}      \\
 {\scriptsize \texttt{del obj[indice]}}      &  {\scriptsize \texttt{obj.\_\_delitem\_\_(indice)}}             \\
 {\scriptsize \texttt{valor in obj}}         &  {\scriptsize \texttt{obj.\_\_contains\_\_(valor)}}             \\
 {\scriptsize \texttt{obj()}}                &  {\scriptsize \texttt{obj.\_\_call\_\_()}}                      \\
 {\scriptsize \texttt{with obj:}}            &  {\scriptsize \texttt{obj.\_\_enter\_\_()}} y {\scriptsize \texttt{obj.\_\_exit\_\_()}}  \\
\end{tabular}
\end{tcolorbox}

\newpage
\subsection{Anexo H: Comparaciones:}

\begin{tcolorbox}
\begin{tabular}{ll}
 \textbf{Operación}  &  \textbf{Método especial}  \\
 {\scriptsize \texttt{a == b}}   &  {\scriptsize \texttt{a.\_\_eq\_\_(b)}}    \\
 {\scriptsize \texttt{a != b}}   &  {\scriptsize \texttt{a.\_\_ne\_\_(b)}}    \\
 {\scriptsize \texttt{a \textless{} b}}    &  {\scriptsize \texttt{a.\_\_lt\_\_(b)}}    \\
 {\scriptsize \texttt{a \textless{}= b}}   &  {\scriptsize \texttt{a.\_\_le\_\_(b)}}    \\
 {\scriptsize \texttt{a \textgreater{} b}}    &  {\scriptsize \texttt{a.\_\_gt\_\_(b)}}    \\
 {\scriptsize \texttt{a \textgreater{}= b}}   &  {\scriptsize \texttt{a.\_\_ge\_\_(b)}}    \\
\end{tabular}
\end{tcolorbox}
\subsection{Anexo I: Operadores aritméticos:}

\begin{tcolorbox}
\begin{tabular}{ll}
 \textbf{Operación}  &  \textbf{Método especial}      \\
 {\scriptsize \texttt{a + b}}    &  {\scriptsize \texttt{a.\_\_add\_\_(b)}}       \\
 {\scriptsize \texttt{a - b}}    &  {\scriptsize \texttt{a.\_\_sub\_\_(b)}}       \\
 {\scriptsize \texttt{a * b}}    &  {\scriptsize \texttt{a.\_\_mul\_\_(b)}}       \\
 {\scriptsize \texttt{a / b}}    &  {\scriptsize \texttt{a.\_\_truediv\_\_(b)}}   \\
 {\scriptsize \texttt{a // b}}   &  {\scriptsize \texttt{a.\_\_floordiv\_\_(b)}}  \\
 {\scriptsize \texttt{a \% b}}    &  {\scriptsize \texttt{a.\_\_mod\_\_(b)}}       \\
 {\scriptsize \texttt{a ** b}}   &  {\scriptsize \texttt{a.\_\_pow\_\_(b)}}       \\
\end{tabular}
\end{tcolorbox}

\newpage
\subsection{Anexo J: Operaciones bit a bit:}

\begin{tcolorbox}
\begin{tabular}{ll}
 \textbf{Operación}  &  \textbf{Método especial}    \\
 {\scriptsize \texttt{a \& b}}    &  {\scriptsize \texttt{a.\_\_and\_\_(b)}}     \\
 {\scriptsize \texttt{a \textbar{} b}}   &  {\scriptsize \texttt{a.\_\_or\_\_(b)}}      \\
 {\scriptsize \texttt{a \string^ b}}    &  {\scriptsize \texttt{a.\_\_xor\_\_(b)}}     \\
 {\scriptsize \texttt{\string~a}}       &  {\scriptsize \texttt{a.\_\_invert\_\_()}}   \\
 {\scriptsize \texttt{a \textless{}\textless{} b}}   &  {\scriptsize \texttt{a.\_\_lshift\_\_(b)}}  \\
 {\scriptsize \texttt{a \textgreater{}\textgreater{} b}}   &  {\scriptsize \texttt{a.\_\_rshift\_\_(b)}}  \\
\end{tabular}
\end{tcolorbox}
\subsection{Anexo K: Gestión de atributos:}

\begin{tcolorbox}
\begin{tabular}{ll}
 \textbf{Operación}                     &  \textbf{Método especial}                  \\
 {\scriptsize \texttt{obj.atributo}} (si no existe)  &  {\scriptsize \texttt{obj.\_\_getattr\_\_(nombre)}}         \\
 {\scriptsize \texttt{obj.atributo}}                 &  {\scriptsize \texttt{obj.\_\_getattribute\_\_(nombre)}}    \\
 {\scriptsize \texttt{obj.atributo = valor}}         &  {\scriptsize \texttt{obj.\_\_setattr\_\_(nombre, valor)}}  \\
 {\scriptsize \texttt{del obj.atributo}}             &  {\scriptsize \texttt{obj.\_\_delattr\_\_(nombre)}}         \\
\end{tabular}
\end{tcolorbox}


\end{document}
